% Supplemental Material: dimension cap, correlation gap, and exceptional gaps.
% Extracted without changes to the proofs from the 16 November 2026 full draft.
\documentclass[11pt,letterpaper,twoside,leqno]{article}
\usepackage[margin=1.05in]{geometry}
\usepackage[T1]{fontenc}
\usepackage{lmodern,microtype}
\usepackage{etoolbox}
\usepackage{amsmath,amssymb,amsthm,mathtools}
\usepackage{longtable,booktabs,tabularx}
\usepackage{xcolor}
\usepackage{tikz}
\usetikzlibrary{arrows.meta}
\usepackage{xr-hyper}
\usepackage[hidelinks,hypertexnames=false,linktoc=all]{hyperref}
\newcommand{\Ac}{\mathsf A}
\usepackage{fancyhdr}
\pagestyle{fancy}
\fancyhf{}
\fancyhead[LE,RO]{\small\thepage}
\fancyhead[LO]{\small\itshape\nouppercase{\leftmark}}
\fancyhead[RE]{\small\itshape Supplemental material: rectangle CLT}
\renewcommand{\headrulewidth}{0.2pt}
\renewcommand{\sectionmark}[1]{\markboth{\thesection.\ #1}{}}
\setlength{\headheight}{14pt}
\hypersetup{pdftitle={Sharp High-Dimensional Gaussian Approximations over Rectangles without Covariance Restrictions},pdfsubject={Independent sums with bounded, sub-exponential, or polynomial-moment envelopes},pdfauthor={Victor Chernozhukov, Denis Chetverikov, Kengo Kato, Yuta Koike}}
\setlength{\parindent}{1.25em}
\setlength{\parskip}{0pt}
\emergencystretch=2em
\allowdisplaybreaks[2]
\raggedbottom
\clubpenalty=10000
\widowpenalty=10000
% Bold statement names and proof headings; italic theorem bodies.
\newtheoremstyle{bookplain}{6pt plus 2pt minus 1pt}{6pt plus 2pt minus 1pt}
  {\itshape}{}{\bfseries}{.}{.5em}
  {\thmname{#1}\thmnumber{ #2}\thmnote{ (#3)}}
\newtheoremstyle{bookremark}{6pt plus 2pt minus 1pt}{6pt plus 2pt minus 1pt}
  {\normalfont}{}{\bfseries}{.}{.5em}
  {\thmname{#1}\thmnumber{ #2}\thmnote{ (#3)}}
\theoremstyle{bookplain}
\newtheorem{theorem}{Theorem}
\newtheorem{proposition}{Proposition}
\newtheorem{lemma}{Lemma}
\newtheorem{corollary}{Corollary}
\newtheorem{conjecture}{Conjecture}
\theoremstyle{bookremark}
\newtheorem{remark}{Remark}
\makeatletter
\renewenvironment{proof}[1][\proofname]{\par
  \pushQED{\qed}\normalfont\topsep6pt plus6pt\relax
  \trivlist\item[\hskip\labelsep\bfseries #1\@addpunct{.}]\ignorespaces
}{\popQED\endtrivlist\@endpefalse}
\makeatother
\newcommand{\E}{\mathbb E}
\newcommand{\Pp}{\mathbb P}
\newcommand{\R}{\mathbb R}
\newcommand{\cR}{\mathcal R}
\newcommand{\cF}{\mathcal F}
\newcommand{\cP}{\mathcal P}
\newcommand{\cD}{\mathcal D}
\newcommand{\Cov}{\operatorname{Cov}}
\newcommand{\Var}{\operatorname{Var}}
\newcommand{\tr}{\operatorname{tr}}
\newcommand{\op}{\mathrm{op}}
\newcommand{\norm}[1]{\left\lVert#1\right\rVert}
\newcommand{\maxnorm}[1]{\lVert#1\rVert_{\max}}
\newcommand{\rh}{\rho_{\mathcal R}}
\newcommand{\ind}{\mathbf 1}
\newcommand{\dd}{\,d}
\newcommand{\Hh}{H_2}

\newcommand{\logp}{\log_+}

\hypersetup{pdftitle={Supplemental Material for Sharp High-Dimensional Gaussian Approximations over Rectangles without Covariance Restrictions},pdfsubject={Dimension-logarithm cap, positive correlation gaps, and exceptional correlation budgets}}
\externaldocument{Rectangle_CLT_main_reference_labels}[Rectangle_CLT_arXiv_2026-11-16.pdf]

\title{\Large Supplemental Material for\\[0.5em]
Sharp High-Dimensional Gaussian Approximations\\[0.15em]
over Rectangles without Covariance Restrictions}
\author{\normalsize Victor Chernozhukov, Denis Chetverikov, Kengo Kato, Yuta
Koike\thanks{This supplement was developed in genuine collaboration
with the AI system GPT-Astra. The human authors take full
responsibility for its content.}}
\date{\normalsize 16 November 2026}

% A compact article title keeps the abstract and contents together.
\makeatletter
\renewcommand{\@maketitle}{%
  \begin{center}
    {\Large\@title\par}\vskip.85em
    {\normalsize\lineskip.5em\begin{tabular}[t]{c}\@author\end{tabular}\par}
    \vskip.65em{\normalsize\@date\par}
  \end{center}}
\makeatother

\begin{document}
\renewcommand{\thepage}{S\arabic{page}}
\maketitle

This supplement to~\cite{MainPaper2026} gives three refinements of
Gaussian approximation over rectangles. Appendix~G caps the additional
logarithm in the bounded-envelope bound. Appendix~H proves the
positive-correlation-gap result in Corollary~\ref{cor:positive-correlation-gap}
of the paper, retaining Gaussian width. Appendix~I proves
Corollary~\ref{cor:correlation-budget}, which replaces a uniform
correlation gap by a budget for exceptional pairs.

We use the notation of the paper throughout. References to
Sections~1--5, Appendices~A--F, and results without a G, H, or I prefix
refer to the paper; references with those prefixes refer to this
supplement. The additional results here are not covered by the
archived Lean certification described in Section~\ref{sec:lean-introduction}
of the paper.

\begingroup
\setcounter{tocdepth}{1}
\small
\tableofcontents
\endgroup
\clearpage
\appendix
\setcounter{section}{6}
\numberwithin{equation}{section}
\numberwithin{theorem}{section}
\numberwithin{lemma}{section}
\numberwithin{proposition}{section}
\numberwithin{corollary}{section}
\numberwithin{remark}{section}

% Dimension cap: adapted from Yuta Koike's 10 November 2026 note, Section 6.
% The coarse single-face bound below avoids importing Source input 3.2.
\section{A dimension-logarithm cap for the bounded-envelope bound}
\label{app:dimension-cap}

The additional logarithm in Theorem~\ref{thm:unrestricted} comes
from integrating a remainder bound proportional to
$(t+\varepsilon)^{-1}$. We can cap this logarithm by
$L=\log(ep)$ by comparing that estimate with two further bounds:
one for terms associated with a single face and another for terms
associated with a pair of faces. These bounds have polynomial
factors in the number of faces and the support rank, but those
factors enter only through logarithms after we integrate the
smaller bound at each time. We use the same interpolation and
cubic estimate as in Section~\ref{sec:smoothing}, retaining the
boundary coefficient from~\eqref{eq:intrinsic-boundary-coefficient}:
\[
 \kappa=\min\{\sqrt L,\,1+\max\{\E\|W\|_\infty,g_\Sigma\}\},
 \qquad 1\le\kappa\le\sqrt L.
\]
This preserves the Gaussian-width improvement when we cap the logarithm.

\begin{proposition}[Dimension-logarithm cap]
\label{prop:dimension-cap}
Under~\eqref{eq:setup}, define $J_L:=\min\{D_L,L\}$. There is an
absolute constant $C$ such that
\begin{equation}
 \rh(W,Z)\le C\min\{1,\delta g_\Sigma L\sqrt{J_L}\}.
 \label{cap:result}
\end{equation}
Here $g_\Sigma=\E\|Z\|_\infty$. More precisely, for
$h\ge\mathsf A\delta$ and $\varepsilon=Lh^2<1/4$, the remainder in
\eqref{eq:path-summary} satisfies
\begin{equation}
 \sup_a\int_0^1|R_{4,a,h}(t)|\,dt
 \le\frac{C\delta^2\kappa L}h
       \min\{1+\log(1/\varepsilon),L\}.
 \label{cap:integrated}
\end{equation}
\end{proposition}

The proposition improves the bound obtained from the envelope
alone. Combining it with~\eqref{eq:intrinsic-unrestricted} gives
\[
 \rh(W,Z)\le C\min\left\{1,
 g_\Sigma L\bigl(\delta+\sqrt{s_4D_L}\bigr),
 \delta g_\Sigma L\sqrt{J_L}\right\}.
\]
Thus the relative fourth-moment bound and the capped envelope
bound can be used together, with the same Gaussian-width factor.
The argument below does not replace $D_L$ by $J_L$ in the expression
involving $s_4$. Using $g_\Sigma\le C\sqrt L$ in~\eqref{cap:result}
gives the dimension-based bound $C\delta L^{3/2}\sqrt{J_L}$.
When $p$ is fixed, both $g_\Sigma$ and $J_L\le L$ are bounded, so
\eqref{cap:result} is of order $\delta$.

We work in the support coordinates of
Subsection~\ref{appx-an:sec:coords}. The rank is $r\le p$,
the unit face normals are $u_1,\ldots,u_m$ with $m=2p$, and
\[
 \sum_i\E\zeta_i^0(\zeta_i^0)^\top=I_r,
 \qquad \|\zeta_i^0\|_U\le\delta,
 \qquad \|x\|_U=\max_{j\le m}|u_j\cdot x|.
\]
We first bound the probabilities of short intervals in one face
projection and small rectangles in two face projections. The
two-dimensional estimate displays explicitly the loss when the
normals are nearly parallel. A later algebraic identity will
compensate for that loss.

\begin{lemma}[Concentration for one or two face projections]
\label{lem:cap-concentration}
Let $V=\sum_iX_i+G_A$, where the $X_i\in\R^r$ are independent
and centered, $\|X_i\|_U\le b\delta$, and $G_A$ is an
independent centered Gaussian vector. Suppose
$\Cov(V)=\tau^2I_r$, with $c_0\le\tau^2\le C_0$.
There are constants $A,C$, depending only on $b,c_0,C_0$,
such that, for $h\ge A\delta$,
\begin{align}
 \sup_x\Pp\{|u_j\cdot V-x|\le h\}&\le Ch,
 \label{cap:one-concentration}\\
 \sup_{x,y}\Pp\{|u_j\cdot V-x|\le h,
                 |u_k\cdot V-y|\le h\}
 &\le Ch^2/\nu_{jk},\qquad
 \nu_{jk}:=\sqrt{1-(u_j\cdot u_k)^2}>0.
 \label{cap:two-concentration}
\end{align}
For the logistic weights $W_j=F\pi_j$ and $W_{jk}=F\pi_j\pi_k$
at width $h$ and arbitrary thresholds, these estimates imply
\begin{equation}
 \sup_z\E W_j(V+z)\le Ch,\qquad
 \sup_z\E W_{jk}(V+z)\le Ch^2/\nu_{jk},
 \label{cap:weight-concentration}
\end{equation}
with the second bound for $\nu_{jk}>0$.
The bounds also hold after an independent random shift.
\end{lemma}

\begin{proof}
We prove the two-dimensional concentration bound first. Fix
$j,k$ with $\nu_{jk}>0$ and write
\[
 \rho:=u_j\cdot u_k,\qquad \nu:=\sqrt{1-\rho^2},\qquad
 q:=\frac{u_k-\rho u_j}{\nu}.
\]
Then $u_j\cdot q=0$ and $\|q\|=1$, so $(u_j,q)$ is an
orthonormal pair. The projected vector
$\widehat V:=(u_j\cdot V,q\cdot V)$ has covariance
$\tau^2I_2$. Its bounded summands
$\widehat X_i:=(u_j\cdot X_i,q\cdot X_i)$ satisfy
\[
 |\widehat X_{i,1}|\le b\delta,
 \qquad
 |\widehat X_{i,2}|
 \le\frac{|u_k\cdot X_i|+|\rho|\,|u_j\cdot X_i|}{\nu}
 \le\frac{2b\delta}{\nu}.
\]
Consequently, on a frequency rectangle
\[
 |t_1|\le c/h,\qquad |t_2|\le c\nu/h,
\]
the scalar phase $t\cdot\widehat X_i$ is bounded in absolute
value by $3cb\delta/h$. Choose the absolute $c$ and then $A$
so that the phase difference for two independent copies is at
most $1$ whenever $h\ge A\delta$.

Let $\widehat X_i'$ be an independent copy. Independence gives
\[
 \left|\E e^{it\cdot\widehat X_i}\right|^2
 =\E\cos\{t\cdot(\widehat X_i-\widehat X_i')\}.
\]
On the chosen frequency rectangle, $1-\cos y\ge c_*y^2$
for these phase differences. Since their variance is
$2\Var(t\cdot\widehat X_i)$, we obtain
\[
 \left|\E e^{it\cdot\widehat X_i}\right|^2
 \le1-2c_*\Var(t\cdot\widehat X_i)
 \le e^{-2c_*\Var(t\cdot\widehat X_i)}.
\]
Multiply over $i$ and include the Gaussian characteristic
function, whose absolute value is the exponential of minus half
its projected variance. The total projected covariance is
$\tau^2I_2$, so the characteristic function of $\widehat V$
satisfies
\[
 |\varphi_{\widehat V}(t_1,t_2)|
 \le e^{-c_*\tau^2(t_1^2+t_2^2)}
\]
throughout that frequency rectangle, after reducing $c_*$ if needed.

We now convert this characteristic-function bound into a
probability bound. The event in~\eqref{cap:two-concentration}
implies
\[
 |\widehat V_1-x|\le h,\qquad
 \left|\widehat V_2-\frac{y-\rho x}{\nu}\right|
 \le\frac{2h}{\nu},
\]
because $u_k\cdot V=\rho\widehat V_1+\nu\widehat V_2$.
Use the nonnegative kernel
$k(z):=(\sin(az)/(az))^2$, with $k(0):=1$ and a fixed
sufficiently small $a>0$. It is bounded below by a positive
constant on $[-1,1]$. Also
$\sin(az)/(az)=(2a)^{-1}\int_{-a}^a e^{itz}\,dt$,
so its square has a bounded triangular Fourier transform
supported on $[-2a,2a]$.
The indicator of the displayed rectangle is therefore bounded
by a constant times
\[
 k\left(\frac{\widehat V_1-x}{h}\right)
 k\left(\frac{\nu\widehat V_2-(y-\rho x)}{2h}\right).
\]
Choose $a$ so that the Fourier support of this product is
inside the frequency rectangle already controlled. Fourier
inversion and the rescaling of the two kernels give
\[
 \begin{split}
 \Pp\{|u_j\cdot V-x|\le h,\ |u_k\cdot V-y|\le h\}
 &\le\frac{Ch^2}{\nu}
   \int_{|t_1|\le c/h,\ |t_2|\le c\nu/h}
       |\varphi_{\widehat V}(t_1,t_2)|\,dt_1dt_2\\
 &\le\frac{Ch^2}{\nu}.
 \end{split}
\]
The last step uses $\tau^2\ge c_0$ and the integral of a
two-dimensional Gaussian exponential over $\R^2$.
The compact Fourier support justifies interchanging the integral
and expectation; no density of $V$ is required. The same argument
with the single projection $u_j\cdot V$ and one rescaled kernel
gives~\eqref{cap:one-concentration}.

For the weight bounds, write
$s=(u_j\cdot x-a_j)/h$ and $t=(u_k\cdot x-a_k)/h$.
Dropping other positive terms from the logistic denominator gives
\[
 W_j(x)\le g_1(s):=\frac{e^s}{(1+e^s)^2},\qquad
 W_{jk}(x)\le g_2(s,t):=
                 \frac{e^{s+t}}{(1+e^s+e^t)^3}.
\]
The changes of variables $u=e^s$ and $w=e^t$ give
\[
 \int_{\R}g_1(s)\,ds=\int_0^\infty\frac{du}{(1+u)^2}=1,
 \qquad
 \int_{\R^2}g_2(s,t)\,ds\,dt
 =\int_0^\infty\int_0^\infty\frac{du\,dw}{(1+u+w)^3}
 =\frac12.
\]
The logarithmic derivative of $g_1$ has absolute value at most
$1$, and both logarithmic partial derivatives of $g_2$ have
absolute value at most $2$. Thus the supremum of either
majorant on a unit interval or square is bounded by an absolute
constant times its integral over that cell. Summing the cell
suprema is therefore finite.

Partition the standardized projection line or plane into these
unit cells. Their inverse images have projection side lengths
$h$, so~\eqref{cap:one-concentration} or
\eqref{cap:two-concentration} bounds the probability of each cell
by $Ch$ or $Ch^2/\nu$. Multiply by the corresponding cell
supremum and sum. This proves~\eqref{cap:weight-concentration}.
A deterministic translation only changes the thresholds, so the
bounds are uniform in that translation. Conditioning on an
independent random shift proves the final assertion.
\end{proof}

We next decompose the Hessian so that the coefficient associated
with two nearly parallel faces becomes visible. For this calculation,
let $F$ have width $\ell$ and define $\mu:=\sum_j\pi_j u_j$.
For fixed directions $z,w$, write $z_j:=u_j\cdot z$,
$w_j:=u_j\cdot w$, and $\rho_{jk}:=u_j\cdot u_k$.
The first derivative of $W_j=F\pi_j$ in direction $u_j$ is
\[
 D_{u_j}W_j=\ell^{-1}W_j(1-2u_j\cdot\mu).
\]
Using this derivative, the Hessian can be written as
\begin{align}
 D^2F(x)[z,w]
 &=-\ell^{-1}\sum_jz_jw_j D_{u_j}W_j(x)
     +2\ell^{-2}\sum_{j<k}W_{jk}(x)B_{jk}(z,w),
 \label{cap:hessian}\\
 B_{jk}(z,w)
 &:=z_jw_k+z_kw_j-\rho_{jk}(z_jw_j+z_kw_k).
 \label{cap:pair-coefficient}
\end{align}
To verify the identity, expand the first term as
$\ell^{-2}\{-\sum_jW_jz_jw_j
          +2\sum_{j,k}\rho_{jk}W_{jk}z_jw_j\}$.
For $j\ne k$, its terms containing $\rho_{jk}$ cancel
the corresponding terms in the unordered-pair sum.
The diagonal terms have $\rho_{jj}=1$, and the remaining terms
are exactly
\[
 \ell^{-2}\left\{2\sum_{j,k}W_{jk}z_jw_k
                         -\sum_jW_jz_jw_j\right\},
\]
the Hessian formula used in Lemma~\ref{appx-an:lem:restoration}.

The fourth-order tensors in~\eqref{appx-an:fourth-identity}
each contain a repeated direction. Recall that $\zeta_i$ and
$\widetilde\zeta_i$ are copies of $\zeta_i^0$, with all copies
mutually independent and independent of the original summands,
masks, and Gaussian vectors. Write the three patterns as
\[
 (Z_i,z_i,w_i)
 =(\zeta_i,\zeta_i,\zeta_i),\qquad
   (\widetilde\zeta_i,\zeta_i,\zeta_i),\qquad
   (\widetilde\zeta_i,\widetilde\zeta_i,\zeta_i).
\]
For each pattern,
$D^4F[Z_i,Z_i,z_i,w_i]=D_{Z_i}^2\{D^2F[z_i,w_i]\}$.
Applying these two derivatives to~\eqref{cap:hessian} gives
terms proportional to
$z_{i,j}w_{i,j}D_{Z_i}^2D_{u_j}W_j$ and
$B_{jk}(z_i,w_i)D_{Z_i}^2W_{jk}$, respectively.
After summing over observations and integrating the scalar
remainder kernels, denote their contributions by
$R_{4,\mathrm{single},a,h}$ and $R_{4,jk,a,h}$.
Thus the first contains the sum over single faces, while the
second is the contribution of the unordered pair $\{j,k\}$.
This specifies the slot order even though the full fourth
derivative is symmetric.

The single-face terms admit a bound with time singularity
$v^{-3/2}$, where $v=t+\varepsilon$. For each pair, we will
prove two bounds: one retains the cancellation as the angle
shrinks, and the other uses the Gaussian part of the path.
Taking their geometric mean removes the angle from the bound.

\begin{lemma}[Complementary remainder bounds]
\label{lem:cap-remainder}
Let $h\ge A\delta$, $\varepsilon=Lh^2<1/4$, and
$v:=t+\varepsilon$. Along the calibrated interpolation,
\begin{align}
 |R_{4,\mathrm{single},a,h}(t)|
 &\le Cmr\delta^2v^{-3/2},
 \label{cap:single}\\
 |R_{4,jk,a,h}(t)|
 &\le C\min\left\{\frac{\delta^2\nu_{jk}}{h^2},
                    \frac{r\delta^2}{\nu_{jk}v}\right\},
 \qquad \nu_{jk}>0.
 \label{cap:pair}
\end{align}
Parallel or antiparallel pairs contribute zero. Consequently,
\begin{equation}
 |R_{4,a,h}(t)|
 \le \frac{Cmr\delta^2}{v^{3/2}}
       +\frac{Cm^2\sqrt r\,\delta^2}{h\sqrt v}.
 \label{cap:complementary}
\end{equation}
\end{lemma}

\begin{proof}
Use the unnormalized path and the rescaled test function in
\eqref{eq:seeded-test}. The logistic width is then
$\ell=\widetilde h=\sqrt{1+\varepsilon}\,h$, comparable to $h$.
As in~\eqref{eq:seeded-gaussian-split}, take independent standard
Gaussian vectors $G_0,G_1$, independent of all summands, masks,
and copies, and write
\[
 \begin{aligned}
 Y_t&=V_{0,t}+\sigma_tG_1,\qquad \sigma_t^2=v/2,\\
 V_{0,t}&=(1+t)\sum_i\mathbb B_{i,t}\zeta_i^0+\sigma_tG_0.
 \end{aligned}
\]
The covariance and summand bounds needed below are
\[
 \Cov(V_{0,t})=(1-t/2+\varepsilon/2)I_r,
 \qquad \|(1+t)\mathbb B_{i,t}\zeta_i^0\|_U\le2\delta.
\]
The covariance scalar lies in $[1/2,9/8]$, so
Lemma~\ref{lem:cap-concentration} applies uniformly in time
to $V_{0,t}$. It also applies to $Y_t$, whose covariance is
$(1+\varepsilon)I_r$. Each remainder point has the form
$Y_t+d_i$ with $\|d_i\|_U\le4\delta$, by
\eqref{appx-an:translates}. We first fix the Taylor parameter
and one of the three copy patterns. The scalar kernels in
\eqref{appx-an:kernel} are common to all indices $i$ and have
uniformly bounded absolute integrals, so we can integrate them
after bounding each pattern.

For the first pair estimate, we make the angular cancellation
in $B_{jk}$ explicit. Fix a nonparallel pair and define
\[
 c:=1-|\rho_{jk}|,\qquad
 s:=\begin{cases}\operatorname{sign}(\rho_{jk}),&\rho_{jk}\ne0,\\
                 1,&\rho_{jk}=0,\end{cases}
 \qquad q:=\frac{u_j-su_k}{\sqrt{2c}}.
\]
The vector $q$ is unit because
$\|u_j-su_k\|^2=2-2|\rho_{jk}|=2c$.
Expanding $(q\cdot z)(q\cdot w)$ gives the identity
\begin{equation}
 B_{jk}(z,w)=sc\{z_jw_j+z_kw_k-2(q\cdot z)(q\cdot w)\}.
 \label{cap:angle-cancellation}
\end{equation}
Every slot in the three patterns has aggregate second-moment
matrix $I_r$. Thus, for any unit vector $e$, Cauchy--Schwarz
over both the observation index and the probability space gives
\[
 \sum_i\E|(e\cdot z_i)(e\cdot w_i)|
 \le\left\{\sum_i\E(e\cdot z_i)^2
             \sum_i\E(e\cdot w_i)^2\right\}^{1/2}=1.
\]
This holds whether the two slots use the same copy or independent
copies. Applying it with $e=u_j,u_k,q$ to
\eqref{cap:angle-cancellation} gives
$\sum_i\E|B_{jk}(z_i,w_i)|\le4c$.
We will also need the pointwise bound
$|B_{jk}(z_i,w_i)|\le4\delta^2$, which follows directly from
\eqref{cap:pair-coefficient} and the face bounds on the copies.

To bound derivatives of the positive weight, differentiate its
logarithm. With $\mu=\sum_l\pi_lu_l$ and
$Z_{i,l}:=u_l\cdot Z_i$, we have
\[
 \begin{aligned}
 D_{Z_i}\log W_{jk}
   &=\ell^{-1}(Z_{i,j}+Z_{i,k}-3\mu\cdot Z_i),\\
 D_{Z_i}^2\log W_{jk}
   &=-3\ell^{-2}\left\{\sum_l\pi_l Z_{i,l}^2
                                 -(\mu\cdot Z_i)^2\right\}.
 \end{aligned}
\]
The braces are a variance for the probabilities $\pi_l$, with
the remaining probability assigned to the value $0$.
Since $\|Z_i\|_U\le\delta$, the two displayed expressions
have absolute values at most $5\delta/\ell$ and
$3\delta^2/\ell^2$. The identity
$D^2W=W\{(D\log W)^2+D^2\log W\}$ therefore yields
\[
 |D_{Z_i}^2W_{jk}(x)|\le C\delta^2\ell^{-2}W_{jk}(x).
\]
The same bound on the first logarithmic derivative shows that
$W_{jk}(Y_t+d_i)\le e^{C\delta/\ell}W_{jk}(Y_t)$.
This factor is bounded because $h\ge A\delta$.
The restored weight at $Y_t$ is independent of all copies.
Combining the preceding bounds, the factor $2\ell^{-2}$ in
\eqref{cap:hessian}, and~\eqref{cap:weight-concentration}, we obtain
\[
 |R_{4,jk,a,h}(t)|
 \le C\delta^2\ell^{-4}c\,\frac{\ell^2}{\nu_{jk}}
 \le\frac{C\delta^2\nu_{jk}}{h^2}.
\]
Here $\nu_{jk}^2=c(1+|\rho_{jk}|)\ge c$, which explains
why the angular cancellation compensates for the concentration
loss in this estimate.

For the other pair estimate and the single-face estimate, we
transfer derivatives to $G_1$ before averaging the copies.
Let $\mathcal H$ be generated by the original summands, masks,
and $G_0$. Then $V_{0,t}$ is $\mathcal H$-measurable, while
$G_1$ is independent of $\mathcal H$. Each repeated direction
is an independent copy of the corresponding original summand,
so
\[
 \sum_i\E[Z_iZ_i^\top\mid\mathcal H]=I_r.
\]
The directions and shifts are jointly independent of $G_1$
conditional on $\mathcal H$; the shifts may still depend on
the copies. This is why simply replacing a shifted weight by
an upper bound inside a signed Gaussian contraction would not
justify the needed averaging.

Define the symmetric Gaussian matrices
\[
 H_2(G):=GG^\top-I_r,
 \qquad
 H_{3,u}(G):=(u\cdot G)(GG^\top-I_r)-uG^\top-Gu^\top
\]
for a unit vector $u$. Conditional on $\mathcal H$ and the
copies, the base, directions, and shift are fixed. Gaussian
integration by parts then gives
\begin{align*}
 \E_G D_Z^2W(V_0+\sigma G+d)
 &=\sigma^{-2}\E_G\!\left[
       \tr\{H_2(G)ZZ^\top\}W(V_0+\sigma G+d)\right],\\
 \E_G D_Z^2D_uW(V_0+\sigma G+d)
 &=\sigma^{-3}\E_G\!\left[
       \tr\{H_{3,u}(G)ZZ^\top\}W(V_0+\sigma G+d)\right].
\end{align*}
In the second identity, the contraction equals
$(u\cdot G)\{(Z\cdot G)^2-\|Z\|^2\}
 -2(u\cdot Z)(Z\cdot G)$: the last two terms arise when the
third integration by parts differentiates the quadratic factor.
The logistic derivatives are bounded for fixed directions, and
the Gaussian polynomial factors are integrable, so these
identities apply to the weights under consideration.

We next remove the copy-dependent shifts by the exact series
\eqref{eq:weight-translation}. This preserves the signed
contraction until the copies have been averaged.
For $k_0=1$ or $2$, write
$W_\alpha:=F\prod_{a=1}^{k_0}\pi_{\alpha_a}$ and define
\[
 b_l(d):=e^{u_l\cdot d/\ell}-1,\qquad
 \eta:=e^{4\delta/\ell}-1.
\]
For our shifts, $|b_l(d_i)|\le\eta$.
Writing the shifted logistic denominator as the unshifted
denominator times $1+\sum_l\pi_l(U)b_l(d)$ gives
\[
 \begin{split}
 W_\alpha(U+d)
 ={}&W_\alpha(U)e^{\sum_{a=1}^{k_0}u_{\alpha_a}\cdot d/\ell}
 \sum_{N=0}^\infty(-1)^N\binom{k_0+N}{N}\\
 &\hspace{7mm}\times
 \sum_{\beta_1,\ldots,\beta_N}
       \prod_{a=1}^N\pi_{\beta_a}(U)b_{\beta_a}(d).
 \end{split}
\]
For $N=0$, the inner sum and product are both $1$.
Choose $A$ large enough that $\eta\le1/2$, so the series
converges absolutely, uniformly over the permitted shifts.

Let $C_i=z_{i,j}w_{i,j}$ for a single-face term and
$C_i=B_{jk}(z_i,w_i)$ for a pair term. In either case
$|C_i|\le4\delta^2$. For a fixed level $N$ and sequence
$\beta=(\beta_1,\ldots,\beta_N)$, the coefficient matrix
after averaging the copies is
\[
 \Xi_{\alpha,\beta}:=
 \sum_i\E\!\left[
 Z_iZ_i^\top C_i
 e^{\sum_{a=1}^{k_0}u_{\alpha_a}\cdot d_i/\ell}
 \prod_{a=1}^N b_{\beta_a}(d_i)
 \,\middle|\,\mathcal H\right].
\]
For any unit vector $x$, the conditional covariance identity gives
\[
 |x^\top\Xi_{\alpha,\beta}x|
 \le4\delta^2e^{4k_0\delta/\ell}\eta^N
       \sum_i\E[(x\cdot Z_i)^2\mid\mathcal H]
 =4\delta^2e^{4k_0\delta/\ell}\eta^N.
\]
Since the matrix is symmetric, this is also a bound on
$\|\Xi_{\alpha,\beta}\|_{\op}$. In particular, the bound
is deterministic even though the matrix may depend on the base.

We can now take absolute values using
$|\tr(H\Xi)|\le\|H\|_*\|\Xi\|_{\op}$, where
$\|\cdot\|_*$ is the nuclear norm. The remaining factors
$\pi_{\beta_a}(U)$ are nonnegative, and their sum over
$\beta$ is $(\sum_l\pi_l(U))^N\le1$.
The further sum over $N$ is bounded because
\[
 \sum_{N=0}^\infty\binom{k_0+N}{N}\eta^N
 =(1-\eta)^{-k_0-1}\le2^{k_0+1}.
\]
Thus the contraction is bounded by a constant times
$\delta^2\|H(G_1)\|_*W_\alpha(V_{0,t}+\sigma_tG_1)$,
apart from the displayed powers of $\sigma_t^{-1}$.
The absolute convergence just used, together with
$\sum_i\E\|Z_i\|^2=r$ and finite Gaussian polynomial
moments, justifies interchanging the series, copy expectations,
and Gaussian expectations.

At this point the coefficient matrices have been removed by
their deterministic bounds. We may therefore average over
the unconditional law of $V_{0,t}$ with $G_1$ fixed.
Independence makes $\sigma_tG_1$ a fixed translation for this
expectation, and~\eqref{cap:weight-concentration} gives bounds
$C\ell$ and $C\ell^2/\nu_{jk}$ for the single and pair
weights. We are using concentration of $V_{0,t}$ itself,
not concentration conditional on $\mathcal H$.

It remains to bound the Gaussian matrix factors. The identities
$\|GG^\top\|_*=\|G\|^2$ and $\|I_r\|_*=r$ give
\[
 \begin{aligned}
 \E\|H_2(G_1)\|_*&\le\E(\|G_1\|^2+r)=2r,\\
 \E\|H_{3,u}(G_1)\|_*
 &\le\E\bigl[|u\cdot G_1|(\|G_1\|^2+r)
                                      +2\|G_1\|\bigr]\le Cr.
 \end{aligned}
\]
For the last inequality, Cauchy--Schwarz bounds
$\E[|u\cdot G_1|\|G_1\|^2]$ by
$\{r(r+2)\}^{1/2}$, while
$\E|u\cdot G_1|\le1$ and $\E\|G_1\|\le\sqrt r$.
Restoring the Hessian coefficients and the Gaussian scales now
yields
\[
 |R_{4,jk,a,h}(t)|
 \le C\ell^{-2}\sigma_t^{-2}\delta^2r
                    \frac{\ell^2}{\nu_{jk}}
 \le\frac{Cr\delta^2}{\nu_{jk}v}.
\]
For each single face the corresponding bound is
$C\ell^{-1}\sigma_t^{-3}\delta^2r\ell
 \le Cr\delta^2v^{-3/2}$.
Summing over the $m$ faces proves~\eqref{cap:single}.
As explained at the start, integration of the scalar kernels
preserves these estimates for the full remainder contributions.

If $|\rho_{jk}|=1$, then $u_k=u_j$ or $u_k=-u_j$,
and substitution in~\eqref{cap:pair-coefficient} gives
$B_{jk}=0$. For every other pair, the smaller of the two
bounds in~\eqref{cap:pair} is at most their geometric mean,
\[
 \left\{\frac{\delta^2\nu_{jk}}{h^2}
         \frac{r\delta^2}{\nu_{jk}v}\right\}^{1/2}
 =\frac{\sqrt r\,\delta^2}{h\sqrt v}.
\]
Summing over at most $m^2/2$ pairs and adding the single-face
bound proves~\eqref{cap:complementary}.
\end{proof}

We finish by integrating the smaller bound at each time.
The two estimates in~\eqref{cap:complementary} have polynomial
factors in $m$ and $r$. Comparing each with the original
$v^{-1}$ estimate confines those factors to logarithms of
the times at which the bounds cross.

\begin{proof}[Proof of Proposition~\ref{prop:dimension-cap}]
Define
\[
 A_h:=\frac{\delta^2\kappa L}h,\qquad
 S_h:=mr\delta^2,\qquad
 B_h:=\frac{m^2\sqrt r\,\delta^2}h.
\]
The envelope allows us to take $s_4=\delta^2$ in
\eqref{eq:seeded-derivatives}. Retaining $\kappa$, that estimate
bounds the full remainder by $CA_h/v$. Lemma~\ref{lem:cap-remainder} bounds the same
remainder by $C(S_h/v^{3/2}+B_h/\sqrt v)$.
Using $\min\{x,y+z\}\le\min\{x,y\}+\min\{x,z\}$ for
nonnegative $x,y,z$, we obtain
\[
 |R_{4,a,h}(t)|\le C\left[
 \min\left\{\frac{A_h}v,\frac{S_h}{v^{3/2}}\right\}
 +\min\left\{\frac{A_h}v,\frac{B_h}{\sqrt v}\right\}\right].
\]

For the first minimum, the two expressions agree at
$v_s:=(S_h/A_h)^2$. Below $v_s$ we use $A_h/v$, and above
$v_s$ we use $S_h/v^{3/2}$. The integral of the latter
over $[v_s,\infty)$ is $2A_h$, while the logarithmic part
is bounded by
$A_h\log_+(v_s/\varepsilon)
 =2A_h\log_+\{S_h/(A_h\sqrt\varepsilon)\}$.
Consequently,
\[
 \int_\varepsilon^{1+\varepsilon}
       \min\{A_h/v,S_h/v^{3/2}\}\,dv
 \le CA_h\left\{1+\log_+
                    \frac{S_h}{A_h\sqrt\varepsilon}\right\}.
\]
If $v_s$ lies outside the integration interval, the same bound
follows by integrating the smaller expression throughout.

For the second minimum the crossing time is
$v_p:=(A_h/B_h)^2$. Below $v_p$ we use $B_h/\sqrt v$,
whose integral over $[0,v_p]$ is $2A_h$; above $v_p$ we use
$A_h/v$. Since $1+\varepsilon<5/4$, this gives
\[
 \int_\varepsilon^{1+\varepsilon}
       \min\{A_h/v,B_h/\sqrt v\}\,dv
 \le CA_h\{1+\log_+(B_h/A_h)\}.
\]
This bound also covers crossing times outside the interval.

The choice $\varepsilon=Lh^2$ removes $h$ from both ratios:
\[
 \frac{S_h}{A_h\sqrt\varepsilon}=\frac{mr}{\kappa L^{3/2}},
 \qquad \frac{B_h}{A_h}=\frac{m^2\sqrt r}{\kappa L}.
\]
Because $\kappa\ge1$, $m=2p$, $r\le p$, and $L=\log(ep)\ge1$,
both positive logarithms are at most $CL$. Thus retaining $\kappa$
does not change the logarithmic cap.
We have proved
$\sup_a\int_0^1|R_{4,a,h}(t)|\,dt\le CA_hL$.
Direct integration of $CA_h/(t+\varepsilon)$ gives the
alternative bound $CA_h\{1+\log(1/\varepsilon)\}$.
Taking their minimum proves~\eqref{cap:integrated}.

We now return to the rectangle discrepancy. The smoothing,
initialization, and cubic estimates in the proof of
Theorem~\ref{thm:unrestricted}, together with the remainder
bound just proved, yield
\[
 \rh(W,Z)\le C\kappa L\left[
 h+\delta+\frac{\delta^2}h
       \min\{1+\log(1/(Lh^2)),L\}\right]
\]
whenever $h\ge\mathsf A\delta$ and $Lh^2<1/4$.
The initialization term $\delta^2/h$ is absorbed here because
the minimum in braces is at least $1$.
Choose $h:=\mathsf A\delta\sqrt{J_L}$, taking the absolute
constant $\mathsf A\ge1$ large enough for the preceding lemmas.
Since $J_L\ge1$, this choice satisfies $h\ge\mathsf A\delta$.
If $Lh^2<1/4$, then $\delta\sqrt L<1$ and
\[
 1+\log\frac1{Lh^2}
 =1+2\log\frac1{\delta\sqrt L}
       -2\log\mathsf A-\log J_L
 \le2D_L.
\]
It follows that the minimum in the discrepancy bound is at
most $2\min\{D_L,L\}=2J_L$.
The three terms in brackets are therefore bounded by
$C\delta\sqrt{J_L}$, giving
$\rh(W,Z)\le C\delta\kappa L\sqrt{J_L}$.

If instead $Lh^2\ge1/4$, then
\[
 \delta\kappa L\sqrt{J_L}
 =\kappa Lh/\mathsf A
 \ge\kappa\sqrt L/(2\mathsf A)\ge1/(2\mathsf A),
\]
so the trivial probability bound gives the same estimate after
increasing the absolute constant. Hence
$\rh(W,Z)\le C\min\{1,\delta\kappa L\sqrt{J_L}\}$ in both cases.

It remains to express the boundary coefficient in terms of Gaussian
width. Lemma~\ref{intrinsic:lem:maximum} gives
$\E\|W\|_\infty\le C(g_\Sigma+\delta L)$ without invoking
the Gaussian approximation being proved. If $\delta L\le g_\Sigma$,
this bound and $g_\Sigma\ge\sqrt{2/\pi}$ imply
$\kappa\le Cg_\Sigma$. If instead $\delta L>g_\Sigma$, then
\[
 \delta g_\Sigma L\sqrt{J_L}\ge g_\Sigma^2\ge2/\pi,
\]
so the trivial probability bound suffices.
This proves~\eqref{cap:result} in all cases.
\end{proof}

\clearpage
\section{Proof of Corollary~\ref*{cor:positive-correlation-gap} (positive correlation gap)}
\label{app:correlation-gap-proof}

We prove the positive-correlation-gap bound using the smoothing and
interpolation estimates established in Section~\ref{sec:smoothing}.
The additional estimate controls the fourth-order remainder after the
Gaussian variance has grown. We retain the intrinsic boundary coefficient
throughout these estimates and express it in terms of $g_\Sigma$ at the end.
The proof uses a cubic Gaussian Hermite
matrix for the single-face terms and retains the angle between the
normals in the two-face terms. Pairing the corresponding surface
contributions gives a bound that remains valid when a three-coordinate
covariance matrix is singular.

\subsection{Reduction to a fourth-order estimate}\label{hgap:sec:reduction}

Throughout this appendix, assume~\eqref{eq:setup} and the positive gap
$c$ in Corollary~\ref{cor:positive-correlation-gap}. Recall the support
coordinates from Section~\ref{sec:smoothing}: $\Sigma=\Gamma\Gamma^\top$
with $\Gamma$ of full column rank $r$, the $m=2p$ unit vectors $u_j$
are its rows and their negatives, and
$\zeta_i^0=\Gamma^\dagger\xi_i$. Thus
\[
 \sum_i\Cov(\zeta_i^0)=I_r,\qquad
 \|\zeta_i^0\|_U\le\delta,
 \qquad \|z\|_U=\max_j|u_j\cdot z|.
\]
For the geometric calculations below, write
$P(a):=\{x\in\R^r:u_j\cdot x\le a_j,\ j\le m\}$, the polyhedron
previously denoted by $A_a$. For real $s$, let $P(a)^s=P(a+s\mathbf1)$.
We allow arbitrary threshold lists, including those defining empty sets.
The logistic function $F_{a,h}$ is the same as in
Section~\ref{sec:smoothing}. We retain the boundary coefficient
from~\eqref{eq:intrinsic-boundary-coefficient},
\[
 \kappa=\min\{\sqrt L,\,1+\max\{\E\|W\|_\infty,g_\Sigma\}\}.
\]
Thus $1\le\kappa\le\sqrt L$. Boundary probabilities contribute the
factor $\kappa$, while Gaussian projection moments supply the remaining
powers of $L$.

Fix $h\ge\mathsf A\delta$ and $\varepsilon:=Lh^2<1/4$, where the
absolute constant $\mathsf A$ will be chosen large enough for the
estimates below. We use the calibrated path~\eqref{eq:path}, namely
\begin{equation}\label{hgap:eq:path}
 U_t=\frac{(1+t)\sum_i\mathbb B_{i,t}\zeta_i^0
                       +\sqrt{t+\varepsilon}\,G}{\sqrt{1+\varepsilon}},
 \qquad \Pp(\mathbb B_{i,t}=1)=\frac{1-t}{(1+t)^2},
 \qquad 0\le t\le1.
\end{equation}
At each time, the masks are independent of one another and of the
summands and standard Gaussian $G$. No joint construction across times
is needed. The normalization makes $\Cov(U_t)=I_r$.

The smoothing error, the short Slepian initialization, and the cubic
term have already been bounded. More precisely, let $I_h$ be the
initialization error, and let $C_{3,a,h}(t)$ and $R_{4,a,h}(t)$ be the
cubic term and fourth-order remainder in the derivative of
$\E F_{a,h}(U_t)$. Lemma~\ref{lem:endpoint} and
Proposition~\ref{an:prop:derivative}, using
$s_3\le\delta$ and $s_4\le\delta^2$, give
\begin{align}
 \rh(W,Z)&\le Ch\kappa L+I_h+
       \sup_a\int_0^1\{|C_{3,a,h}(t)|+|R_{4,a,h}(t)|\}\dd t,
                                             \label{hgap:eq:reduction}\\
 I_h+\sup_a\int_0^1|C_{3,a,h}(t)|\dd t
     &\le C\kappa L(\delta+\delta^2/h),        \label{hgap:eq:old-cubic}\\
 |R_{4,a,h}(t)|&\le\frac{C\delta^2\kappa L}{h(t+\varepsilon)}.
                                             \label{hgap:eq:early}
\end{align}
All three estimates are uniform over the thresholds. The substitutions
for $s_3,s_4$ follow directly by bounding each remaining face coordinate
by $\delta$ after the repeated pair has been summed using its covariance.

The only additional interpolation estimate needed is
\begin{equation}\label{hgap:eq:late}
 |R_{4,a,h}(t)|\le
 \frac{C\delta^2\kappa L^{3/2}}{d_c(t+\varepsilon)^{3/2}},
 \qquad d_c:=c(2-c).
\end{equation}
We prove it in Subsection~\ref{hgap:sec:completion}, using the Gaussian
estimates developed next. Taking the smaller of
\eqref{hgap:eq:early} and~\eqref{hgap:eq:late} before integration will
replace the endpoint logarithm by $1+\log(1/c)$. This is the only
place where the proof of the present corollary requires more than
the estimates used for Theorem~\ref{thm:unrestricted}.

\subsection{Gaussian surface estimates}\label{hgap:sec:surfaces}

The additional derivative estimate will be applied at a sum of two
independent vectors: a base $V$ that supplies boundary concentration,
and a Gaussian vector that supplies integration by parts. We first
state the property of the base that is needed throughout the argument.
As in~\eqref{appx-an:concentration}, we call $V$ a
\emph{concentrated base} if, for a fixed absolute $C_0$,
\begin{equation}\label{hgap:eq:ac}
 \Pp\{V\in P(a)^s\setminus P(a)^{-s}\}
       \le C_0\kappa(s+\delta),
       \qquad a\in\R^m,\ s>0.
\end{equation}
The coefficient $\kappa$ is fixed throughout, and all constants in
the following estimates are independent of it. The same inequality holds
for a subsystem of the faces by sending
unused thresholds to infinity. It also survives an independent random
translation, by conditioning on that translation. Throughout the
additional estimates, $G\sim N(0,I_r)$ is independent of $V$ and
$U:=V+\sigma G$. We use
\(\delta\sqrt L\le\sigma\le1\) and a smoothing width \(h\ge\delta\),
except that the surface estimates below also apply when
\(0<\sigma<\delta\sqrt L\). Write \([m]=\{1,\ldots,m\}\) for the
indices of the \(m=2p\) signed face normals.

Fix a threshold list $a$. We use the logistic weights already introduced
in~\eqref{appx-an:basic-notation}, which we recall here to specify their
dependence on the evaluation point:
\[
 q_j(x):=\frac{u_j\cdot x-a_j}{h},\qquad
 F(x):=\frac1{1+\sum_j e^{q_j(x)}},\qquad
 \pi_j(x):=e^{q_j(x)}F(x),\qquad \theta(x):=1-F(x).
\]
For an ordered tuple $\alpha=(\alpha_1,\ldots,\alpha_k)\in[m]^k$, define
\[
 W_\alpha(x):=F(x)\prod_{b=1}^k\pi_{\alpha_b}(x).
\]
Repeated indices are allowed. These products arise when a translated
logistic weight is expanded at a common point. Their sums are simple:
$\sum_{\alpha\in[m]^k}W_\alpha=F\theta^k\le F\theta$ for $k\ge1$.
The notation $\norm{\cdot}_*$ denotes the nuclear norm of a matrix.
Each weight $W_\alpha$ is log-concave: its logarithm is an affine
function minus $(k+1)\log(1+\sum_j e^{q_j})$.
The correlation gap $c>0$ means
$|u_j\cdot u_k|\le1-c$ for normals from distinct original coordinates.
The antipodal normals of a single coordinate are excepted. For
linearly independent normals, write
\[
 \vartheta_{jk}:=\arccos(u_j\cdot u_k),\qquad
 \lambda_{jk}:=\vartheta_{jk}\sin\vartheta_{jk},\qquad
 d_c=c(2-c),
\]
and let $\cD$ be the set of their ordered index pairs. Then
\begin{equation}
 \lambda_{jk}\ge\sin^2\vartheta_{jk}\ge d_c
 \qquad((j,k)\in\cD).
 \label{hgap:eq:gap}
\end{equation}

We next prove the geometric estimates that will control the signed
traces of the Gaussian matrices. A \emph{facet} is a face of codimension
$1$, and a \emph{ridge} is a face of codimension $2$. The estimates use
two cancellations among their boundary contributions.
Pairing the boundary contributions from neighboring facets preserves
cancellation even after the facet thresholds are weighted by their
absolute values. At the next codimension, weighting
each ridge by its exterior angle produces a spherical divergence identity.
Both estimates can therefore be localized to short outward tubes.

Let $\varphi_r$ be the standard Gaussian density on $\R^r$, and let
$\mathcal H^q$ denote $q$-dimensional Hausdorff measure.
For a proper full-dimensional polyhedron $P$ with the specified normals,
write $F_j$ for its actual facets, with supporting equation
$u_j\cdot g=b_j$. Let $E$ run over its faces of codimension $2$.
Such a face is shared by exactly two facets, and $\vartheta_E$ denotes
the angle between their outward normals. Define
\begin{align}
 J_{\mathrm f}(P)&:=\sum_j |b_j|
   \int_{F_j}(\norm g^2-r)\varphi_r(g)\dd\mathcal H^{r-1}(g),
 \label{hgap:eq:Jf}\\
 J_{\mathrm e}(P)&:=\sum_E \vartheta_E
   \int_E(\norm g^2-r)\varphi_r(g)\dd\mathcal H^{r-2}(g).
 \label{hgap:eq:Je}
\end{align}
Empty sets and the whole space have both quantities equal to $0$.
There are no codimension-$2$ faces when $r=1$.

\begin{lemma}[Two localized surface estimates]\label{hgap:lem:surfaces}
Let $V$ satisfy~\eqref{hgap:eq:ac}, and let $Q=P(a)$ be empty or full
dimensional. For every $0<\sigma\le1$, define
$(Q-V)/\sigma:=\{g:V+\sigma g\in Q\}$. Then
\begin{equation}
 \E\bigl|J_{\mathrm f}((Q-V)/\sigma)\bigr|
 +\E\bigl|J_{\mathrm e}((Q-V)/\sigma)\bigr|
 \le C\kappa\bigl(\sigma L^{3/2}+\delta L^2\bigr).
 \label{hgap:eq:surface-bounds}
\end{equation}
The constant depends only on $C_0$, and no angle separation is needed.
\end{lemma}

\begin{proof}
Let $R(g)$ be $1$ plus the largest length of the projection of $g$
onto a span of at most three linearly independent face normals.
Each projection norm is $1$-Lipschitz, and so is their maximum.
There are at most $m+m^2+m^3$ relevant subspaces. For a subspace $S$
of dimension at most $3$, the projected Gaussian has the law of a
standard Gaussian in that dimension, and
$\E\exp(\|\Pi_SG\|^2/4)\le2^{3/2}$.
Markov's inequality and a union bound give
\[
 \Pp\{R(G)>1+x\}\le
       \min\{1,2^{3/2}(m+m^2+m^3)e^{-x^2/4}\}.
\]
Integrating this tail and using $\log m\le CL$ proves
\begin{equation}
 \E R(G)^q\le C_qL^{q/2},\qquad q=1,2,3,4.
 \label{hgap:eq:R-moments}
\end{equation}
The reason for using spans of up to three normals is that the
integration-by-parts formulas below involve faces of codimension at
most $3$. The projection of a point onto any such normal space has
norm at most $R$ at that point. We will bound each absolute surface
quantity by
\begin{equation}
 C\E_G\left[R(G)^4
       \ind\{G\in P^{1/R(G)}\setminus P\}\right].
 \label{hgap:eq:localized}
\end{equation}

For the facet quantity, define $s_j:=\int_{F_j}\varphi_r$ and
$\rho_{jk}:=\int_{E_{jk}}\varphi_r$.
We give the Gauss--Green calculation explicitly. On $F_j$, write
$g=b_ju_j+y$ with $y\perp u_j$. The tangential divergence is
\[
 \operatorname{div}_{F_j}(y\varphi_r(g))
       =(r-1-\|y\|^2)\varphi_r(g).
\]
Since $\|g\|^2-r=(b_j^2-1)+(\|y\|^2-(r-1))$, integration over
$F_j$ and summation of the boundary terms give
\[
 J_{\mathrm f}(P)
 =\sum_j |b_j|(b_j^2-1)s_j
 -\sum_{\{j,k\}}
   \bigl(|b_j|d_{k\mid j}+|b_k|d_{j\mid k}\bigr)\rho_{jk},
\]
where, for $\rho:=u_j\cdot u_k$,
\[
 n_{k\mid j}:=\frac{u_k-\rho u_j}{\sin\vartheta_{jk}},\qquad
 d_{k\mid j}:=\frac{b_k-\rho b_j}{\sin\vartheta_{jk}}
             =g\cdot n_{k\mid j}\quad(g\in E_{jk}).
\]
Here $n_{k\mid j}$ is the outward conormal within $F_j$.
On the ridge, let $z$ be the projection onto
$\operatorname{span}(u_j,u_k)$. Then
\[
 |d_{k\mid j}|\le\norm z,\quad
 |d_{k\mid j}+d_{j\mid k}|\le\vartheta_{jk}\norm z,\quad
 \bigl||b_j|-|b_k|\bigr|\le\vartheta_{jk}\norm z.
\]
Indeed, the sum of the two conormals has norm
$2\sin(\vartheta_{jk}/2)\le\vartheta_{jk}$, and
$|b_j-b_k|=|(u_j-u_k)\cdot z|\le\vartheta_{jk}\|z\|$.
These estimates control the paired term without dividing by a small
angle.
To use this cancellation, write the paired coefficient as
\[
 |b_j|(d_{k\mid j}+d_{j\mid k})
                  +( |b_k|-|b_j| )d_{j\mid k}.
\]
Since $|b_j|\le\|z\|$, the three preceding inequalities give
\begin{equation}
 \bigl||b_j|d_{k\mid j}+|b_k|d_{j\mid k}\bigr|
 \le 2\vartheta_{jk}\norm z^2.
 \label{hgap:eq:weighted-pair}
\end{equation}
On a facet the coefficient $|b_j|(b_j^2-1)$ in the
Gauss--Green formula has absolute value at most $CR(x)^3$.

For the ridge quantity, let $z_E$ be the projection onto the
two-dimensional normal space of $E$, and write $\rho_E:=\int_E\varphi_r$.
The same divergence calculation in its $(r-2)$-dimensional tangent
space now has $\|g\|^2-r=(\|z_E\|^2-2)
 +(\|g-z_E\|^2-(r-2))$. Thus Gauss--Green within each ridge gives
\begin{equation}
 J_{\mathrm e}(P)
 =\sum_E\vartheta_E(\norm{z_E}^2-2)\rho_E
 -\sum_H\left\langle z_H,
        \sum_{E\supset H}\vartheta_E n_{E,H}\right\rangle\rho_H.
 \label{hgap:eq:edge-identity}
\end{equation}
Here $H$ runs over codimension-$3$ faces,
$\rho_H:=\int_H\varphi_r\dd\mathcal H^{r-3}$, and $z_H$ is the
projection of a point of $H$ onto its three-dimensional normal space.
This projection is constant on $H$. The vector $n_{E,H}$ is the
outward unit conormal to $H$ within $E$: it is tangent to the affine
hull of $E$, orthogonal to $H$, and points out of $E$ at $H$.
The second sum is absent when $r\le2$.

It remains to control the conormals in this second sum. Fix $H$ and
work in its three-dimensional normal space. The normal cone
$C_H$ consists of nonnegative combinations of normals to the facets
containing $H$. Its intersection $\Omega_H:=C_H\cap S^2$ is a spherical
polygon. A ridge $E$ containing $H$ corresponds to an edge of this
polygon, joining the two normals of the facets meeting at $E$.
The edge has length $\vartheta_E$, and its outward spherical conormal
is $-n_{E,H}$. Indeed, the remaining normals generating $C_H$ have
nonnegative inner product with $n_{E,H}$, so the outward direction
from that edge is the negative one.

For a constant vector $v$ in this normal space, define the tangent
field $X_v(\omega):=v-(v\cdot\omega)\omega$ on $S^2$.
The function $v\cdot\omega$ has spherical Laplacian
$-2v\cdot\omega$, and hence
$\operatorname{div}_{S^2}X_v=-2v\cdot\omega$.
On the edge corresponding to $E$, its flux is
$X_v\cdot(-n_{E,H})=-v\cdot n_{E,H}$, because
$\omega\cdot n_{E,H}=0$. The spherical divergence theorem gives
\[
 -\sum_{E\supset H}\vartheta_E v\cdot n_{E,H}
 =-2\int_{\Omega_H}v\cdot\omega\dd\mathcal H^2(\omega).
\]
This holds for every $v$, and therefore
\begin{equation}
 \sum_{E\supset H}\vartheta_E n_{E,H}
 =2\int_{\Omega_H}\omega\dd\mathcal H^2(\omega).
 \label{hgap:eq:spherical-balance}
\end{equation}
In particular, its norm is at most $2|\Omega_H|$. This identity holds
for a spherical polygon with any number of edges, so no simplicity
assumption on the codimension-$3$ face is required.

We now convert these surface bounds into a Gaussian probability of
an outer layer. For a face $F$ of codimension $q\in\{1,2,3\}$, let
$N_F$ be its outward normal cone. For $x\in\operatorname{relint}F$,
use the points
\[
 g=x+y,\qquad y\in N_F,\qquad 0<\|y\|<\frac1{4R(x)}.
\]
We verify disjointness before adding the integrals over these regions.
Every $y\in N_F$ is a nonnegative combination of normals active at
$x$. Thus $y\cdot(w-x)\le0$ for every $w\in P$, and
\[
 \|g-w\|^2=\|y\|^2+\|x-w\|^2+2y\cdot(x-w)
              \ge\|y\|^2+\|x-w\|^2.
\]
It follows that $x$ is the unique nearest point to $g$ in $P$.
The relative interiors of distinct faces are disjoint. Hence a point
cannot belong to regions based on two different faces, since its
nearest point would have to belong to both relative interiors.

The tangent space of $F$ is orthogonal to the span of $N_F$, so the
volume element in these coordinates is
$\dd\mathcal H^{r-q}(x)\dd\mathcal H^q(y)$.
The normal space is spanned by $q$ linearly independent active
normals, which are among those used to define $R$. Consequently,
$|x\cdot y|\le R(x)\|y\|\le1/4$, and
\[
 \frac{\varphi_r(x+y)}{\varphi_r(x)}
       =\exp\{-x\cdot y-\|y\|^2/2\}
\]
is bounded above and below by absolute positive constants. The
$1$-Lipschitz property and $R(x)\ge1$ give the same comparability for
$R(x+y)$ and $R(x)$.

The volume of the truncated normal cone is
$1/(4R(x))$ for a facet,
$\vartheta_E/(32R(x)^2)$ for a ridge, and
$|\Omega_H|/(192R(x)^3)$ for a codimension-$3$ face.
An integral of $R(g)^4\varphi_r(g)$ over the corresponding region
therefore bounds, up to an absolute factor, the surface integrals with
integrands
\[
 R(x)^3\varphi_r(x),\qquad
 \vartheta_E R(x)^2\varphi_r(x),\qquad
 |\Omega_H|R(x)\varphi_r(x),
\]
respectively. These are exactly the bounds needed for the facet term,
the paired ridge terms in~\eqref{hgap:eq:weighted-pair}, and the
codimension-$3$ terms in~\eqref{hgap:eq:edge-identity}, using
\eqref{hgap:eq:spherical-balance}.

Every such point $g=x+y$ lies outside $P$. Also,
$u_j\cdot g\le b_j+\|y\|$ for every $j$, while
$R(g)\le R(x)+\|y\|$ implies $\|y\|R(g)<1$.
Thus all the regions lie in $P^{1/R(g)}\setminus P$.
Their disjointness proves~\eqref{hgap:eq:localized} for both surface
quantities. For unbounded faces, one may first multiply the tangent
fields by smooth cutoffs equal to $1$ on a ball of radius $M$.
The cutoff derivatives are supported where $\|g\|\ge M$, and their
integrals tend to zero by Gaussian integrability of polynomials.
Letting $M$ increase proves the same Gauss--Green identities without
a boundedness assumption on the faces.

Finally, for $P=(Q-V)/\sigma$, the layer event implies
$V+\sigma G\in Q^{\sigma/R(G)}\setminus Q$.
Conditioning on $G$ and applying~\eqref{hgap:eq:ac} bounds
the average in~\eqref{hgap:eq:localized} by
\[
 C\kappa\{\sigma\E R(G)^3+\delta\E R(G)^4\}
 \le C\kappa(\sigma L^{3/2}+\delta L^2).
\]
The projection-moment estimates retain their powers of $L$; only the
boundary probability supplies $\kappa$. This proves both assertions.
\end{proof}

We will also use the positive versions of the same tube estimates.
They follow by applying the tube construction directly to the indicated
nonnegative integrands, without integration by parts. With $P=(Q-V)/\sigma$
and $\sigma\ge\delta\sqrt L$, they give
\begin{align}
 \E\sum_E\vartheta_E\rho_E&\le C\sigma\kappa\sqrt L,
 &\E\sum_E\vartheta_E\norm{z_E}^2\rho_E&\le C\sigma\kappa L^{3/2}.
 \label{hgap:eq:positive-edge}
\end{align}
The first uses $R^2$ on the tubes, giving
$C\kappa(\sigma\E R(G)+\delta\E R(G)^2)\le C\sigma\kappa\sqrt L$.
The second uses $R^4$ and $\|z_E\|^2\le R^2$, giving the second
bound in exactly the same way.

\subsection{Gaussian matrices and logistic weights}\label{hgap:sec:matrices}

To use the surface estimates for smooth derivatives, we first express
powers of the logistic function as averages of polyhedron indicators.
This also explains the scale factors when a threshold is differentiated.
For an integer
$q\ge1$, let $T_q$ be Gamma with shape $q$ and rate $1$, and let
$\omega_1,\ldots,\omega_m$ be independent exponentials of mean $1$,
independent of $T_q,V,G$. Conditional on $T_q$, the probability that all exponentials exceed
$T_qe^{q_j(x)}$ is $\exp\{-T_q\sum_j e^{q_j(x)}\}$. Integrating the
Gamma density therefore gives the identity
\begin{equation}
 F(x)^q=\Pp\{u_j\cdot x\le a_j+h\log(\omega_j/T_q),\ j\le m\}.
 \label{hgap:eq:mixture}
\end{equation}
Here the probability is over the auxiliary Gamma and exponential
variables with $x$ fixed. When $x=V+\sigma G$, we first condition on
$V,G$ and then use their independence from these auxiliary variables.
For a realization of the latter, denote the resulting polyhedron in
$x$-coordinates by $Q$ and set $P:=(Q-V)/\sigma$.
Its $j$th threshold in $g$-coordinates is
$(a_j+h\log(\omega_j/T_q)-u_j\cdot V)/\sigma$.

For a polynomially bounded weight $\psi$ independent of $a$, define
$\mathcal I_\psi(P):=\int_P\psi(g)\varphi_r(g)\dd g$.
Integration first in the normal coordinate of a facet gives
\begin{equation}\label{hgap:eq:coarea-one}
 \partial_{a_j}\mathcal I_\psi(P)
     =\frac1\sigma\int_{F_j(P)}\psi(g)\varphi_r(g)
                                  \dd\mathcal H^{r-1}(g).
\end{equation}
For independent $u_j,u_k$, the Gram determinant of the map
$g\mapsto(u_j\cdot g,u_k\cdot g)$ is $\sin^2\vartheta_{jk}$.
Integration in these two normal coordinates therefore gives
\begin{equation}\label{hgap:eq:coarea-two}
 \partial_{a_j}\partial_{a_k}\mathcal I_\psi(P)
     =\frac1{\sigma^2\sin\vartheta_{jk}}
       \int_{E_{jk}(P)}\psi(g)\varphi_r(g)
                                  \dd\mathcal H^{r-2}(g).
\end{equation}
A nonexistent facet or ridge contributes zero. The formulas hold at
almost every threshold list. To justify averaging them, first use a
threshold difference quotient, or a rectangular difference quotient
for two thresholds. Its absolute value is bounded
by the supremum of the corresponding Gaussian slice integral with the
absolute weight. At fixed rank this supremum is finite for every weight
used below because the Gaussian density in the normal coordinates dominates
its polynomial growth, and the remaining Gaussian coordinates have
finite moments. The factors from the linear change of variables are
$1/\sigma$ and $1/(\sigma^2\sin\vartheta_{jk})$, respectively.
Conditional on $T_q$, the independent threshold shifts have densities,
so the quotients converge at almost every shifted threshold. The
exceptional lists at which a nonempty polyhedron has empty interior,
or a redundant hyperplane contains an entire facet or ridge, have
Lebesgue measure zero. Every such degeneracy forces a linear relation
among some thresholds whose normals are linearly dependent; there
are only finitely many subsets of normals. Consequently the surface
estimates apply almost surely to the mixture polyhedra. Dominated
convergence permits averaging the slice formulas. Only linearly
independent pairs enter the mixed identity.

We also record the elementary consequence of~\eqref{hgap:eq:ac} that
will control the logistic masses. Write $F_s:=F_{a+s\mathbf1,h}$.
The one-variable logistic calculation in
Appendix~\ref{appx-sec:smoothing} gives
$F\theta\le (e/2)(F_h-F_{-h})$. Apply~\eqref{hgap:eq:mixture} with
$q=1$ to the difference and then use~\eqref{hgap:eq:ac} at the randomly
shifted thresholds. Since $h\ge\delta$, this gives
\begin{equation}\label{hgap:eq:logistic-mass}
 \sup_z\E_V F(V+z)\theta(V+z)\le Ch\kappa.
\end{equation}
Using just face $j$ gives the same bound for
$e^{q_j(V+z)}/(1+e^{q_j(V+z)})^2$. The shifts $z$ here are arbitrary
deterministic vectors, so these inequalities may be used after fixing
$G$. All concentration statements concern the unconditional law of $V$.

The matrix estimates have a common reduction to tilted means and
signed traces. We record it separately so that the two uses below
have the same normalization. The reference weight in the mean bound
is chosen to depend only on the face coordinates already present in
the logistic weight; this prevents a logarithm of the smoothing width.

\begin{lemma}[Gaussian matrices and projected reference weights]
\label{hgap:lem:gaussian}
Let $G\sim N(0,I_r)$, and let $w\ge0$ have positive Gaussian mass
$p_w:=\E w(G)$ and finite Gaussian second moments. Define the weighted
probability law and its mean and covariance by
\[
 \dd P_w(g):=\frac{w(g)\varphi_r(g)}{p_w}\dd g,
 \qquad \mu_w:=\E_{P_w}g,\qquad K_w:=\Cov_{P_w}(g).
\]
If $w$ is log-concave and $w(g)\le C_w(1+\|g\|)$, then
$K_w\preceq I_r$ and
\begin{equation}\label{hgap:eq:matrix-reduction}
 \left\|\E[(GG^\top-I_r)w(G)]\right\|_*
 \le 2p_w\|\mu_w\|^2-\E[(\|G\|^2-r)w(G)].
\end{equation}
If instead, or in addition, $w\le b$, where $b$ depends only on an
orthogonal projection $\Pi g$ and $p_b:=\E b(G)<\infty$, then
\begin{equation}\label{hgap:eq:relative-mean}
 p_w\|\mu_w\|^2
 \le\E[w(G)\|\Pi G\|^2]+2p_w\log(p_b/p_w).
\end{equation}
Zero-mass weights may be omitted from either inequality.
\end{lemma}

\begin{proof}
We first establish covariance contraction, following the argument in
Appendix~\ref{appx-an:sec:matrix}. Define
$f(x):=\int w(g)\varphi_r(g-x)\dd g$.
Pr\'ekopa's theorem~\cite{Prekopa1973} makes $f$ log-concave.
The growth bound on $w$ and Gaussian exponential integrability justify
differentiation locally in $x$. At zero the derivatives are
\[
 f(0)=p_w,\qquad \nabla f(0)=\E[Gw(G)],\qquad
 \nabla^2 f(0)=\E[(GG^\top-I_r)w(G)].
\]
It follows that $\nabla^2\log f(0)=K_w-I_r\preceq0$.
Both matrices $p_w\mu_w\mu_w^\top$ and $p_w(I_r-K_w)$ are positive
semidefinite, and the nuclear norm of such a matrix equals its trace.
Writing
\[
 M:=\E[(GG^\top-I_r)w(G)]
       =p_w\mu_w\mu_w^\top-p_w(I_r-K_w),
\]
we obtain
$\|M\|_*\le p_w\|\mu_w\|^2+p_w\tr(I_r-K_w)
=2p_w\|\mu_w\|^2-\tr M$.
Since $\tr M=\E[(\|G\|^2-r)w(G)]$, this proves the first inequality.

For the second inequality, separate the components of $\mu_w$ in
and orthogonal to the range of $\Pi$. For $v$ orthogonal to that range,
Jensen's inequality under $P_w$ and $w\le b$ give
\[
 p_we^{v\cdot\mu_w}
 \le p_w\E_{P_w}e^{v\cdot g}
 =\E[w(G)e^{v\cdot G}]
 \le\E[b(G)e^{v\cdot G}]
 =p_be^{\|v\|^2/2}.
\]
The last equality uses independence of $\Pi G$ and $v\cdot G$.
Choose $v=(I_r-\Pi)\mu_w$, so that $v\cdot\mu_w=\|v\|^2$.
Taking logarithms yields
$\|(I_r-\Pi)\mu_w\|^2\le2\log(p_b/p_w)$.
For the other component, Cauchy--Schwarz gives
\[
 p_w\|\Pi\mu_w\|^2
 =\frac{\|\E[w(G)\Pi G]\|^2}{p_w}
 \le\E[w(G)\|\Pi G\|^2].
\]
Adding these bounds proves~\eqref{hgap:eq:relative-mean}.
\end{proof}

For later use, write $H_2(g):=\norm g^2-r$ and introduce the
\emph{cubic Hermite matrix in direction $u$}
\begin{equation}
 \mathcal H_{3,u}(g):=(u\cdot g)(gg^\top-I_r)-ug^\top-gu^\top.
 \label{hgap:eq:H3}
\end{equation}
The next result includes arbitrary additional face factors. This
uniformity is needed to restore derivatives evaluated at translated bases.

\begin{proposition}[Single-face cubic matrix estimate]\label{hgap:prop:cubic}
Suppose $V$ satisfies~\eqref{hgap:eq:ac}, $G\sim N(0,I_r)$ is
independent of $V$, and $U=V+\sigma G$, with
$h\ge\delta$ and $\delta\sqrt L\le\sigma\le1$.
For every integer $k\ge1$,
\begin{equation}
 \E_V\sum_{\alpha\in[m]^k}
 \left\|\E_G\bigl[\mathcal H_{3,u_{\alpha_1}}(G)
                    W_\alpha(U)\bigr]\right\|_*
 \le C^{k+1}h\kappa L^{3/2}.
 \label{hgap:eq:cubic-matrix}
\end{equation}
No correlation gap is needed for this estimate.
\end{proposition}

\begin{proof}
Fix $V$ while taking Gaussian expectations. For a tuple whose first
index is $j$, we first bound the part
$\E_G[(u_j\cdot G)(GG^\top-I_r)W_\alpha(U)]$.
Split $(u_j\cdot g)W_\alpha(V+\sigma g)$ into the difference of the weights
\[
 w_{\alpha,+}(g):=(u_j\cdot g)_+W_\alpha(V+\sigma g),\qquad
 w_{\alpha,-}(g):=(-u_j\cdot g)_+W_\alpha(V+\sigma g).
\]
Each weight is log-concave: its linear factor is positive on a
halfspace, its logarithm is concave there, and $W_\alpha$ is log-concave.
It is bounded by $\|g\|$, so Lemma~\ref{hgap:lem:gaussian} applies.
Apply~\eqref{hgap:eq:matrix-reduction} to the two weights and use the
triangle inequality. Their trace terms add to
$-\E[H_2(G)|u_j\cdot G|W_\alpha(U)]$.

We first bound the mean terms in this application. With $V$ fixed,
define
\[
 b_j(g):=\frac{e^{q_j(V+\sigma g)}}{(1+e^{q_j(V+\sigma g)})^2},
 \qquad b_{\alpha,\pm}(g):=(\pm u_j\cdot g)_+b_j(g).
\]
The inequality $W_\alpha\le W_j\le b_j$ shows that these reference
weights dominate $w_{\alpha,\pm}$. They depend only on $u_j\cdot g$.
The reference bounds the tilted mean through
\eqref{hgap:eq:relative-mean}, with $\Pi$ the projection onto the line
spanned by $u_j$. We write $p_{\alpha,\pm},\mu_{\alpha,\pm}$ for the
mass and mean associated with $w_{\alpha,\pm}$, and
$p_\alpha,\mu_\alpha$ for those associated with $W_\alpha$.
All of these quantities may depend on $V$.

Use~\eqref{hgap:eq:logistic-mass} to average the weights over $V$.
Define $M_G:=\max_j|u_j\cdot G|$. A union bound for scalar Gaussian
tails gives $\E M_G^q\le C_qL^{q/2}$ for $q=1,2,3$.
After summing the tuples and the two signs, the projection term in
\eqref{hgap:eq:relative-mean} is bounded by
\[
 \E\sum_\alpha |u_{\alpha_1}\cdot G|^3W_\alpha(U)
 \le\E[M_G^3 F(U)\theta(U)]
 \le Ch\kappa\,\E M_G^3\le Ch\kappa L^{3/2}.
\]
In the penultimate step we condition on $G$ and apply
\eqref{hgap:eq:logistic-mass} to the deterministic shift $\sigma G$.
The same argument with $M_G$ instead of $M_G^3$ bounds the total
averaged mass of $w_{\alpha,\pm}$ by $Ch\kappa\sqrt L$.
For each first index $j$, the sum of the two reference weights is
$|u_j\cdot G|b_j(G)$. The single-face version of
\eqref{hgap:eq:logistic-mass} gives
$\E[|u_j\cdot G|b_j(G)]\le Ch\kappa\,\E|u_j\cdot G|
\le Ch\kappa$. There are $m^{k-1}$ tuples with first index $j$,
so their total reference mass is at most $Cm^kh\kappa$.
To see the entropy bound without introducing $\log(1/h)$, let $M$
and $B$ be these total averaged masses. If $M>0$, the weights
$p_{\alpha,\pm}/M$, together with the law of $V$, define a probability
measure on the joint space of $V$, tuples, and signs. Jensen's inequality
for the concave logarithm gives
\[
 \frac1M\E_V\sum_{\alpha,\pm}p_{\alpha,\pm}
       \log\frac{p_{b_{\alpha,\pm}}}{p_{\alpha,\pm}}
 \le\log\left\{\frac1M\E_V\sum_{\alpha,\pm}
                                      p_{b_{\alpha,\pm}}\right\}
 =\log(B/M).
\]
This is the log-sum inequality with the outer average included.
If $M=0$, all terms vanish. Write $M=h\kappa\sqrt L\,x$, where
$0\le x\le C$. Since $B\le Cm^kh\kappa$ and $\log m\le CL$, the
inequalities $x\log(1/x)\le1/e$ and $L\ge1$ give
$M\log(B/M)\le C(k+1)h\kappa L^{3/2}$. The common factor
$h\kappa$ cancels inside the logarithm, so this step introduces no
additional dependence on the concentration coefficient. The zero-mass case is interpreted
by continuity. This proves
\begin{equation}
 \E_V\sum_{\alpha,\pm}p_{\alpha,\pm}
                   \norm{\mu_{\alpha,\pm}}^2\le C(k+1)h\kappa L^{3/2}.
 \label{hgap:eq:cubic-means}
\end{equation}

For the trace, summing tuples leaves
$\sum_j|u_j\cdot G|F\pi_j\theta^{k-1}$.
Use
$F\pi_j\theta^{k-1}=-(h/k)\partial_{a_j}\theta^k$ and expand
$\theta^k=\sum_{q=0}^k(-1)^q\binom{k}{q}F^q$.
For each $q\ge1$, formula~\eqref{hgap:eq:coarea-one}, with
$\psi(g)=(\|g\|^2-r)|u_j\cdot g|$, places $g$ on the $j$th facet.
There $|u_j\cdot g|$ equals the absolute value of its threshold in
$g$-coordinates. Hence the sum over $j$ is precisely
$\sigma^{-1}J_{\mathrm f}((Q-V)/\sigma)$, averaged over the mixing
variables for $F^q$. The $q=0$ term has zero derivative, and the sum
of the remaining absolute coefficients is at most $2^kh/(k\sigma)$.
By Lemma~\ref{hgap:lem:surfaces} and
$\delta\sqrt L\le\sigma$, the averaged absolute surface integral is
at most $C\sigma\kappa L^{3/2}$. We conclude that
\begin{equation}
 \left|\E H_2(G)\sum_j|u_j\cdot G|F(U)\pi_j(U)\theta(U)^{k-1}\right|
 \le C2^kh\kappa L^{3/2}.
 \label{hgap:eq:cubic-trace}
\end{equation}
Finally, the two remaining terms of~\eqref{hgap:eq:H3} have nuclear norm
at most $2\norm{\E_G[G W_\alpha(U)]}$.
For these terms apply~\eqref{hgap:eq:relative-mean} to $W_\alpha$
with the reference $b_j$. The total mass is at most $Ch\kappa$.
The summed projection term is bounded by
$Ch\kappa\,\E M_G^2\le Ch\kappa L$, and the reference mass is at
most $Cm^kh\kappa$. The same log-sum argument, now writing the
total mass as $h\kappa$ times a bounded scalar, gives
\[
 \E_V\sum_\alpha p_\alpha\le Ch\kappa,\qquad
 \E_V\sum_\alpha p_\alpha\norm{\mu_\alpha}^2
 \le C(k+1)h\kappa L.
\]
Since $\E_G[GW_\alpha]=p_\alpha\mu_\alpha$, Cauchy--Schwarz on the
joint space of $V$ and tuples gives
\[
 \E_V\sum_\alpha p_\alpha\|\mu_\alpha\|
 \le\left(\E_V\sum_\alpha p_\alpha\right)^{1/2}
      \left(\E_V\sum_\alpha p_\alpha\|\mu_\alpha\|^2\right)^{1/2}
 \le C\sqrt{k+1}\,h\kappa\sqrt L.
\]
Together with~\eqref{hgap:eq:cubic-means}--\eqref{hgap:eq:cubic-trace}, this
proves~\eqref{hgap:eq:cubic-matrix}.
\end{proof}

For the two-face term, the angle must be retained until after the
matrix inequality is applied. Its role is to make the trace become
$J_{\mathrm e}$, whose codimension-$3$ contributions cancel by
\eqref{hgap:eq:spherical-balance}.

\begin{proposition}[Angle-weighted two-face matrix estimate]\label{hgap:prop:pair}
Suppose $V$ satisfies~\eqref{hgap:eq:ac}, $G\sim N(0,I_r)$ is
independent of $V$, and $U=V+\sigma G$, with
$h\ge\delta$ and $\delta\sqrt L\le\sigma\le1$.
For every integer $K\ge2$,
\begin{equation}
 \E_V\sum_{\substack{\alpha\in[m]^K\\(\alpha_1,\alpha_2)\in\cD}}
 \lambda_{\alpha_1\alpha_2}
 \left\|\E_G[(GG^\top-I_r)W_\alpha(U)]\right\|_*
 \le C^{K+1}\frac{h^2\kappa L^{3/2}}{\sigma}.
 \label{hgap:eq:weighted-pair-matrix}
\end{equation}
In particular, under the correlation gap the same sum without the
weights $\lambda$ is at most $C^{K+1}h^2\kappa L^{3/2}/(d_c\sigma)$.
\end{proposition}

\begin{proof}
Apply~\eqref{hgap:eq:matrix-reduction} to each $W_\alpha$, and multiply
by $\lambda_{jk}$ for $j=\alpha_1,k=\alpha_2$.
To bound its mean, take $\Pi$ to be the projection onto
$\operatorname{span}(u_j,u_k)$, denoted by $\Pi_{jk}$,
in~\eqref{hgap:eq:relative-mean}, with the reference weight
\[
 b_{jk}(g):=\frac{e^{q_j(V+\sigma g)+q_k(V+\sigma g)}}
                  {(1+e^{q_j(V+\sigma g)}+e^{q_k(V+\sigma g)})^3}
                  \ge W_\alpha(V+\sigma g).
\]
Writing the two scores as $s,t$ and substituting $x=e^s,y=e^t$
gives their integral explicitly:
\[
 \int_{\R^2}\frac{e^{s+t}}{(1+e^s+e^t)^3}\dd s\dd t
 =\int_0^\infty\int_0^\infty\frac{\dd x\dd y}{(1+x+y)^3}
 =\frac12.
\]
We next bound the mass of this reference weight. Write
$\rho:=u_j\cdot u_k$ and
$e:=(u_k-\rho u_j)/\sin\vartheta_{jk}$.
The variables $G_j:=u_j\cdot G$ and $G_e:=e\cdot G$ are independent
standard Gaussians, independent of $V$, and
\[
 u_j\cdot U=u_j\cdot V+\sigma G_j,\qquad
 u_k\cdot U=u_k\cdot V+\sigma\rho G_j
                           +\sigma\sin\vartheta_{jk}G_e.
\]
Conditional on $V,G_j$, the second coordinate has density at most
$C/(\sigma\sin\vartheta_{jk})$. The density of the first coordinate
is at most $C\kappa$. To verify this last assertion, partition the
possible values of $u_j\cdot V$ into intervals of length $\sigma$
around a proposed density argument. The marginal form of
\eqref{hgap:eq:ac} bounds each interval probability by
$C\kappa(\sigma+\delta)$. On the interval at distance comparable
to $|\ell|\sigma$, the Gaussian kernel is at most
$C\sigma^{-1}e^{-c_1\ell^2}$. Summing gives
$C\kappa(1+\delta/\sigma)\le C\kappa$.

More explicitly, the joint density at $(x,y)$ is the expectation over
$V$ of the product of the first Gaussian density at $x$ and the
conditional second Gaussian density at $y$. Bounding the latter and
averaging the former therefore gives
\[
 f_{j,k}(x,y)\le\frac{C\kappa}{\sigma\sin\vartheta_{jk}}.
\]
Changing from the two coordinates to their scores contributes $h^2$.
The score integral above and
$\lambda_{jk}=\vartheta_{jk}\sin\vartheta_{jk}$ thus give
$\lambda_{jk}\E b_{jk}(G)\le Ch^2\kappa/\sigma$.
Counting the at most $m^K$ tuples yields total weighted reference
mass at most $Cm^Kh^2\kappa/\sigma$.

For the actual logistic mass we can sum the faces before estimating.
The identity $\partial_{a_j}\partial_{a_k}F=2h^{-2}W_{jk}$,
the mixture~\eqref{hgap:eq:mixture}, and
\eqref{hgap:eq:coarea-two} show that
\[
 \E\sum_{(j,k)\in\cD}\lambda_{jk}W_{jk}(U)
 =\frac{h^2}{\sigma^2}
       \E_{V,\omega,T_1}\sum_E\vartheta_E\rho_E.
\]
The factor $2$ from the ordered pairs cancels the $1/2$ in the
threshold derivative. With the additional weight
$\|\Pi_{jk}G\|^2$, its restriction to the ridge equals
$\|z_E\|^2$. The two positive surface estimates
\eqref{hgap:eq:positive-edge} therefore give
\begin{align}
 \E\sum_{(j,k)\in\cD}\lambda_{jk}W_{jk}(U)
   &\le C h^2\kappa\sqrt L/\sigma,\label{hgap:eq:pair-mass}\\
 \E\sum_{(j,k)\in\cD}\lambda_{jk}W_{jk}(U)\norm{\Pi_{jk}G}^2
   &\le C h^2\kappa L^{3/2}/\sigma.\label{hgap:eq:pair-projection}
\end{align}
Counting the two orders of each pair changes only the absolute constant.
The factor $1/\sin\vartheta_{jk}$ in coarea leaves exactly
$\vartheta_{jk}$. Additional face factors sum to $\theta^{K-2}\le1$,
so the same upper bounds apply after summing all tuples.
For clarity, multiply both masses by the same angle weight before
applying the log-sum inequality. If
$M:=\E_V\sum_\alpha\lambda_{\alpha_1\alpha_2}p_\alpha$ and
$B:=\E_V\sum_\alpha\lambda_{\alpha_1\alpha_2}p_{b_{\alpha_1\alpha_2}}$,
where the sums are over tuples with independent first two normals, then
\[
 \E_V\sum_\alpha\lambda_{\alpha_1\alpha_2}p_\alpha
       \log\frac{p_{b_{\alpha_1\alpha_2}}}{p_\alpha}
 \le M\log(B/M).
\]
Here $M\le Ch^2\kappa\sqrt L/\sigma$ and
$B\le Cm^Kh^2\kappa/\sigma$. Writing
$M=(h^2\kappa\sqrt L/\sigma)x$ shows that the right side is at most
$C(K+1)h^2\kappa L^{3/2}/\sigma$, including the zero-mass case by
continuity. The common factor $h^2\kappa/\sigma$ cancels inside the
logarithm. Thus no logarithm of an angle, of $h$, or of $\kappa$ enters
this step.
Together with~\eqref{hgap:eq:pair-projection}, this gives
\begin{equation}
 \E_V\sum_{\alpha:\,(\alpha_1,\alpha_2)\in\cD}
 \lambda_{\alpha_1\alpha_2}p_\alpha\norm{\mu_\alpha}^2
 \le C(K+1)h^2\kappa L^{3/2}/\sigma.
 \label{hgap:eq:pair-means}
\end{equation}

It remains to control the summed trace. Summing the extra indices
leaves the weight
\[
 \sum_{(j,k)\in\cD}\lambda_{jk}F\pi_j\pi_k\theta^{K-2}.
\]
Expand
$\theta^{K-2}=\sum_{q=0}^{K-2}(-1)^q\binom{K-2}{q}F^q$ and use
\[
 F^{q+1}\pi_j\pi_k
 =\frac{h^2}{(q+1)(q+2)}
       \partial_{a_j}\partial_{a_k}F^{q+1}.
\]
Apply~\eqref{hgap:eq:coarea-two} to each power $F^{q+1}$ with
$\psi(g)=\|g\|^2-r$. Its factor
$(\sigma^2\sin\vartheta_{jk})^{-1}$, multiplied by
$\lambda_{jk}$, leaves $\vartheta_{jk}/\sigma^2$.
After summing the two orders of every ridge, the trace is a signed
sum of mixtures of $J_{\mathrm e}((Q-V)/\sigma)$.
The absolute coefficients sum to at most $C2^Kh^2/\sigma^2$.
By Lemma~\ref{hgap:lem:surfaces}, its magnitude is at most
$C2^K h^2\kappa L^{3/2}/\sigma$. Combining this with~\eqref{hgap:eq:pair-means}
and~\eqref{hgap:eq:matrix-reduction} proves~\eqref{hgap:eq:weighted-pair-matrix}.
The unweighted conclusion follows from~\eqref{hgap:eq:gap}.
\end{proof}

\subsection{Fourth derivatives at translated bases}\label{hgap:sec:translated}

The matrix bounds above were proved at a common point $V+\sigma G$.
The interpolation remainder instead evaluates different summands at
different nearby points. We first decompose the remaining Hessian
into the single-face and two-face terms just bounded, and then expand
the translated weights so that all summands can be contracted at the
common point.
The algebraic decomposition retains cancellation within the Hessian.
For $z_j:=u_j\cdot z$, $w_j:=u_j\cdot w$, and
$\rho_{jk}:=u_j\cdot u_k$,
\begin{equation}
 D^2F(x)[z,w]
 =-\frac1h\sum_j z_jw_jD_{u_j}W_j(x)
 +\frac2{h^2}\sum_{j\ne k}W_{jk}(x)z_j(w_k-\rho_{jk}w_j).
 \label{hgap:eq:hessian}
\end{equation}
To verify the identity, differentiate $D F[w]=-h^{-1}\sum_jw_jW_j$
to obtain
\[
 D^2F[z,w]=-h^{-2}\sum_jz_jw_jW_j
                    +2h^{-2}\sum_{j,k}z_jw_kW_{jk}.
\]
Also,
$D_{u_j}W_j=h^{-1}W_j(1-2\sum_k\rho_{jk}\pi_k)$.
Substitution of this formula into the first term of
\eqref{hgap:eq:hessian} leaves exactly the second displayed sum,
since the terms containing $\rho_{jk}z_jw_jW_{jk}$ cancel.
The summand with $j=k$ in the two-face correction is zero.
If $u_k=\pm u_j$, the corresponding coefficient in the second sum is
$0$. Thus only pairs in $\cD$ contribute.

\begin{lemma}[Fourth derivatives at translated concentrated bases]
\label{hgap:thm:fourth}
Let $\mathcal H$ be a sigma-field, let $V$ be an $\mathcal H$-measurable
concentrated base satisfying~\eqref{hgap:eq:ac}, and let
$G\sim N(0,I_r)$ be independent of $\mathcal H$.
Suppose the finite array $(Z_i,z_i,w_i,v_i)_{i=1}^n$, with all entries
in $\R^r$, is jointly independent of $G$ conditional on $\mathcal H$.
For a fixed $b\ge0$, suppose almost surely that
\[
 \sum_i\E[Z_iZ_i^\top\mid\mathcal H]\preceq I_r,\qquad
 \norm{z_i}_U,\norm{w_i}_U\le\delta,\qquad
 \norm{v_i}_U\le b\delta.
\]
If $\delta\sqrt L\le\sigma\le1$, there are constants $A_b\ge1,C_b$
such that $h\ge A_b\delta$ implies
\begin{equation}
 \left|\sum_i\E D^4F(V+\sigma G+v_i)[Z_i,Z_i,z_i,w_i]\right|
 \le \frac{C_b}{d_c}\frac{\delta^2\kappa L^{3/2}}{\sigma^3}.
 \label{hgap:eq:fourth-bound}
\end{equation}
The constants depend only on $b,C_0$, and not on the rank. No mutual independence of the
entries of an array is required, and $v_i$ may depend on the other
entries. Concentration~\eqref{hgap:eq:ac} is required only unconditionally,
not after conditioning on $\mathcal H$.
\end{lemma}

\begin{proof}
First condition on $\mathcal H$ and on all entries of the arrays.
By the independence assumption, $G$ remains standard Gaussian under
this conditioning, and the vectors used as derivative directions are
fixed. Two integrations by parts in the direction $Z_i$ give
\[
 \E_G D^4F(U+v_i)[Z_i,Z_i,z_i,w_i]
 =\frac1{\sigma^2}\E_G\left[
       \tr\{(GG^\top-I_r)Z_iZ_i^\top\}
                       D^2F(U+v_i)[z_i,w_i]\right].
\]
We insert~\eqref{hgap:eq:hessian} into this expression. For its first term, integrate once more
in the unit direction $u_j$ to obtain
\[
 \E_G[(GG^\top-I_r)D_{u_j}W_j(U+v_i)]
 =\sigma^{-1}\E_G[\mathcal H_{3,u_j}(G)W_j(U+v_i)].
\]
The scalar face coefficients satisfy
$|z_{i,j}w_{i,j}|\le\delta^2$ and
$|z_{i,j}(w_{i,k}-\rho_{jk}w_{i,j})|\le2\delta^2$.
We must nevertheless retain the translation $v_i$ until we can
separate the Gaussian matrix from the summand-dependent coefficient.
Bounding each summand individually at this point would lose the
aggregate second-moment bound.

For this separation, expand the translated weights in the small
changes of their exponential scores. Define
$d_j(v):=e^{u_j\cdot v/h}-1$ and
$\eta_b:=e^{b\delta/h}-1$. For a weight with $K=1$ or $2$ face factors,
\begin{align*}
 W_\alpha(x+v)
 &=W_\alpha(x)e^{\sum_{a=1}^K u_{\alpha_a}\cdot v/h}
     \bigl(1+\sum_j\pi_j(x)d_j(v)\bigr)^{-K-1}\\
 &=e^{\sum_{a=1}^K u_{\alpha_a}\cdot v/h}
   \sum_{\ell\ge0}(-1)^\ell\binom{K+\ell}{\ell}
   \sum_{\beta\in[m]^\ell}W_{\alpha,\beta}(x)
                       \prod_{a=1}^\ell d_{\beta_a}(v).
\end{align*}
We now average the arrays conditional on $\mathcal H$ while keeping
$G$ fixed. Each unshifted weight depends only on $V,G$, so it can be
separated from the summand-dependent coefficient. More explicitly,
for $\alpha=(j)$ or $\alpha=(j,k)$, define
\[
 a_{i,\alpha}:=
 \begin{cases}
 z_{i,j}w_{i,j},&K=1,\\
 z_{i,j}(w_{i,k}-\rho_{jk}w_{i,j}),&K=2.
 \end{cases}
\]
Apart from the deterministic binomial coefficient and its sign, the
matrix multiplying $W_{\alpha,\beta}(U)$ is
\[
 M_{\alpha,\beta}:=\sum_i\E\left[
 a_{i,\alpha}\exp\!\left\{\sum_{a=1}^K u_{\alpha_a}\cdot v_i/h\right\}
 \prod_{a=1}^\ell d_{\beta_a}(v_i)\,Z_iZ_i^\top
 \,\middle|\,\mathcal H\right].
\]
The absolute value of its scalar coefficient is at most
$2\delta^2e^{Kb\delta/h}\eta_b^\ell$.
For any scalar coefficients bounded by $M_0$, the aggregate
second-moment assumption gives, in every unit direction $e$,
\[
 \left|\sum_i\E[\text{coefficient}_i(e\cdot Z_i)^2
                       \mid\mathcal H]\right|
 \le M_0\sum_i\E[(e\cdot Z_i)^2\mid\mathcal H]\le M_0.
\]
Thus $\|M_{\alpha,\beta}\|_{\op}\le
2\delta^2e^{Kb\delta/h}\eta_b^\ell$.
Their dependence on $\mathcal H$ causes no difficulty: operator--nuclear
norm duality bounds their contraction pointwise, after which the matrix
estimates are averaged over the unconditional law of $V$.
Apply Proposition~\ref{hgap:prop:cubic} when $K=1$, and
Proposition~\ref{hgap:prop:pair} when $K=2$.
Choose $A_b$ so that $C\eta_b\le1/2$.
For $K=1,2$, the resulting series is bounded because
\[
 \sum_{\ell\ge0}\binom{K+\ell}{\ell}(C\eta_b)^\ell
       =(1-C\eta_b)^{-K-1}\le2^{K+1}.
\]
The single-face contribution has the factor $h^{-1}\sigma^{-3}$
from the Hessian decomposition and the three integrations by parts.
Proposition~\ref{hgap:prop:cubic} supplies $h\kappa L^{3/2}$ for its matrix sum.
The two-face contribution has factor $h^{-2}\sigma^{-2}$, and
Proposition~\ref{hgap:prop:pair} supplies $h^2\kappa L^{3/2}/(d_c\sigma)$.
Consequently their sum is bounded by
\[
 C_b\delta^2\left\{\frac{h\kappa L^{3/2}}{h\sigma^3}
             +\frac{h^2\kappa L^{3/2}}{h^2\sigma^2d_c\sigma}\right\}
 \le\frac{C_b\delta^2\kappa L^{3/2}}{d_c\sigma^3},
\]
where we used $d_c\le1$. This is~\eqref{hgap:eq:fourth-bound}.
For $\eta_b<1$ the weight expansion converges absolutely since
$\sum_j\pi_j\le1$. At fixed rank, $\sum_i\E\|Z_i\|^2\le r$ and
Gaussian moments through degree three justify the initial interchanges.
The estimates above then remove rank from the final bound.
\end{proof}


\subsection{Application to the interpolation and completion of the proof}
\label{hgap:sec:completion}

We now verify that this translated-derivative bound applies to the
remainder in~\eqref{hgap:eq:path}. This step is essential because the
remainder points depend on the same summands that appear as derivative
directions. The conditional formulation above was chosen to allow
exactly that dependence.

\begin{proof}[Proof of the late bound~\eqref{hgap:eq:late}]
Fix $0<t<1$ and set $Y_t:=\sqrt{1+\varepsilon}\,U_t$. With
$\widetilde a:=\sqrt{1+\varepsilon}\,a$ and
$\widetilde h:=\sqrt{1+\varepsilon}\,h$, we have
$F_{a,h}(U_t)=f(Y_t)$ for $f:=F_{\widetilde a,\widetilde h}$.
Take independent standard Gaussians $G_0,G_1$, independent also of
the summands and masks, and write
\begin{equation}\label{hgap:eq:split}
 Y_t=V_{0,t}+\sigma_tG_1,\qquad
 V_{0,t}:=(1+t)\sum_i\mathbb B_{i,t}\zeta_i^0+\sigma_tG_0,
 \qquad \sigma_t^2:=\frac{t+\varepsilon}{2}.
\end{equation}
The retained sum and $\sigma_tG_0$ are independent and have common
face variances $1-t$ and $(t+\varepsilon)/2$, respectively. At least one
variance is at least $1/3$. The retained summands are centered and
have face norm at most $2\delta$. By~\eqref{appx-an:retained-maximum},
the retained sum has expected face norm at most $2\E\|W\|_\infty$.
The Gaussian component has expected face norm $\sigma_tg_\Sigma\le g_\Sigma$.
Apply both the intrinsic and dimension-based parts of
Lemma~\ref{lem:concentration} to the component with variance at least
$1/3$, treating the other component as an independent shift. The two
bounds give~\eqref{hgap:eq:ac} with the coefficient $\kappa$ recalled
above, with an absolute $C_0$ uniformly in $t$. Also, if $h\ge\sqrt2\delta$, then
\[
 \sigma_t^2\ge Lh^2/2\ge\delta^2L,
 \qquad \sigma_t^2<(1+1/4)/2<1.
\]
Thus all scale conditions are satisfied.

Let $\mathcal H$ be generated by the original summands, masks, and
$G_0$. Let $(\zeta_i)_i$ and $(\widetilde\zeta_i)_i$ be independent
copies of the original array, independent also of $\mathcal H,G_1$.
For clarity, Lemma~\ref{appx-an:lem:identity} writes
its fourth-order remainder as
\begin{equation}\label{hgap:eq:explicit-remainder}
\begin{split}
 R_{4,a,h}(t)=\int_0^1\sum_i\E\bigl[\,
 &a_t(s)D^4f(Y_t+v_{i,t,s})[\zeta_i,\zeta_i,\zeta_i,\zeta_i]\\
 {}+{}&b_t(s)D^4f(Y_t+v_{i,t,s})
                    [\widetilde\zeta_i,\widetilde\zeta_i,\zeta_i,\zeta_i]\\
 {}+{}&c_t(s)D^4f(Y_t+v_{i,t,s})
           [\widetilde\zeta_i,\widetilde\zeta_i,\widetilde\zeta_i,\zeta_i]
                              \,\bigr]\dd s,
\end{split}
\end{equation}
where
\[
 v_{i,t,s}:=(1+t)(s\zeta_i-\mathbb B_{i,t}\zeta_i^0),
 \qquad \|v_{i,t,s}\|_U\le4\delta.
\]
For $\eta_t:=(1-t)/(1+t)$, its scalar kernels are
\begin{align*}
 a_t(s)&:=\frac{(1+t)^2}{6}(1-s)^2
                  \{\eta_t-1+(1+2\eta_t)s\},\\
 b_t(s)&:=\frac{1-t}{2}(1-s),\qquad
 c_t(s):=\frac{1-t}{6}(1-\eta_t).
\end{align*}
They are common to all $i$ and satisfy
$\int_0^1(|a_t|+|b_t|+|c_t|)\dd s\le C$ uniformly in $t$.

In each of the three terms in~\eqref{hgap:eq:explicit-remainder}, the first
two slots form a repeated copy $Z_i$ with
$\sum_i\E[Z_iZ_i^\top\mid\mathcal H]=I_r$.
Each remaining slot has face norm at most $\delta$. The copies and
shifts are jointly independent of $G_1$ conditional on $\mathcal H$.
Lemma~\ref{hgap:thm:fourth}, with $b=4$, therefore bounds each summed
expectation by $C\delta^2\kappa L^{3/2}/(d_c\sigma_t^3)$ whenever
$h\ge A_4\delta$. Since $\widetilde h\ge h$, the test function $f$
has sufficient width. Integration against the scalar kernels gives
\[
 |R_{4,a,h}(t)|\le\frac{C\delta^2\kappa L^{3/2}}{d_c\sigma_t^3}
     \le\frac{C\delta^2\kappa L^{3/2}}{d_c(t+\varepsilon)^{3/2}},
\]
which is~\eqref{hgap:eq:late}.
\end{proof}

The remaining step is a one-dimensional integration. We include it
because it explains why a power of the gap in a pointwise derivative
bound produces only a logarithm in the final error.

\begin{proof}[Proof of Corollary~\ref{cor:positive-correlation-gap}]
Define $J_c:=1+\log(1/c)$. First suppose $h\ge\mathsf A\delta$ and
$\varepsilon=Lh^2<1/4$, where $\mathsf A$ is large enough for all preceding estimates.
Combining~\eqref{hgap:eq:early} and~\eqref{hgap:eq:late}, and then setting
$y:=(t+\varepsilon)/\varepsilon$ and using
$\sqrt\varepsilon=h\sqrt L$, gives
\begin{align}
 \int_0^1|R_{4,a,h}(t)|\dd t
 &\le\frac{C\delta^2\kappa L}h
    \int_1^\infty\min\{y^{-1},d_c^{-1}y^{-3/2}\}\dd y\notag\\
 &=\frac{C\delta^2\kappa L}h
                      \{2\log(1/d_c)+2\}
 \le\frac{C\delta^2\kappa LJ_c}{h}.              \label{hgap:eq:time-integral}
\end{align}
Indeed, the two integrands agree at $y=d_c^{-2}$, the first integrates
to $2\log(1/d_c)$ before that point, and the second integrates to $2$
after it. The extension of the upper integration limit to infinity
is an upper bound, so no assumption is needed about whether the
crossover occurs before the end of the actual interpolation.
Both remainder bounds have the same factor $\kappa$, so it does not
affect the crossover. Finally, $c\le d_c\le1$ gives the last inequality.

Substituting~\eqref{hgap:eq:old-cubic} and~\eqref{hgap:eq:time-integral} into
\eqref{hgap:eq:reduction}, and using $J_c\ge1$, yields
\begin{equation}\label{hgap:eq:balance}
 \rh(W,Z)\le C\kappa L\left(h+\delta+\frac{\delta^2J_c}{h}\right).
\end{equation}
Choose $h:=\mathsf A\delta\sqrt{J_c}$. If $Lh^2<1/4$, the right
side of~\eqref{hgap:eq:balance} is at most $C\delta\kappa L\sqrt{J_c}$.
If $Lh^2\ge1/4$, then
\[
 \delta\kappa L\sqrt{J_c}
 =\frac{\kappa Lh}{\mathsf A}
 \ge\frac{\kappa\sqrt L}{2\mathsf A}\ge\frac1{2\mathsf A},
\]
because $\kappa,L\ge1$. The bound $\rh(W,Z)\le1$ therefore gives
\[
 \rh(W,Z)\le C\min\{1,\delta\kappa L\sqrt{J_c}\}
\]
in both cases.

It remains to express the boundary coefficient in terms of Gaussian
width. Lemma~\ref{intrinsic:lem:maximum} gives
$\E\|W\|_\infty\le C(g_\Sigma+\delta L)$ without using the rectangle
approximation being proved. If $\delta L\le g_\Sigma$, this estimate
and $g_\Sigma\ge\sqrt{2/\pi}$ imply $\kappa\le Cg_\Sigma$.
If instead $\delta L>g_\Sigma$, then
\[
 \delta g_\Sigma L\sqrt{J_c}\ge g_\Sigma^2\ge2/\pi,
\]
so the trivial bound suffices. Hence
$\rh(W,Z)\le C\min\{1,\delta g_\Sigma L\sqrt{J_c}\}$ for all
admissible parameters, as asserted.
\end{proof}

The proof also covers $p=1$: there are no independent face pairs, the
two-face term vanishes, and $c=d_c=J_c=1$. More generally, singular
covariances cause no loss because every estimate takes place on the
Gaussian support. In particular, the normal-cone identity
\eqref{hgap:eq:spherical-balance} imposes no three-coordinate
nondegeneracy condition.


\clearpage
\section{Proof of Corollary~\ref*{cor:correlation-budget} (exceptional correlation gaps)}
\label{app:correlation-budget}

The correlation budget allows pairs to have arbitrarily small
positive gaps by charging their aggregate contribution. We therefore use two estimates for the fourth-order
remainder. Lemma~\ref{lem:cap-remainder} bounds each exceptional pair
at a cost that vanishes with its gap. The matrix estimates in
Appendix~\ref{app:correlation-gap-proof} control the remaining pairs
together and retain the boundary coefficient $\kappa$. We separate
the two contributions and explain why the matrix estimates apply to
the separated pairs, before choosing a smoothing width that
accommodates both contributions.

\begin{proof}[Proof of Corollary~\ref{cor:correlation-budget}]
Fix $c\in(0,1]$, and recall
\[
 \eta_{jk}:=1-|\Sigma_{jk}|,\qquad
 B_c:=\sum_{\substack{j<k\\0<\eta_{jk}<c}}\sqrt{\eta_{jk}},
 \qquad J_c:=1+\log(1/c).
\]
Here $j,k$ index the original $p$ coordinates. For the derivative
estimates we work in the Gaussian support, with the $m=2p$ signed
unit face normals from Section~\ref{sec:smoothing}. Each coordinate
provides a normal and its negative. We retain the coefficient
\[
 \kappa:=\min\{\sqrt L,\,1+\max\{\E\|W\|_\infty,g_\Sigma\}\},
 \qquad 1\le\kappa\le\sqrt L,
\]
until the last step, where we replace it by the Gaussian width.
Initially take $h\ge\mathsf A\delta$ and $\varepsilon:=Lh^2<1/4$,
where $\mathsf A$ is an absolute constant large enough for the
estimates in Appendices~\ref{app:dimension-cap}
and~\ref{app:correlation-gap-proof}. Write $v:=t+\varepsilon$.

We use the calibrated interpolation from those appendices. To check
the hypotheses needed below, recall its Gaussian split
\eqref{hgap:eq:split}:
\[
 Y_t=V_{0,t}+\sigma_tG_1,\qquad
 V_{0,t}=(1+t)\sum_i\mathbb B_{i,t}\zeta_i^0+\sigma_tG_0,
 \qquad \sigma_t^2=v/2.
\]
The standard Gaussians $G_0,G_1$ are independent of one another and
of the original summands and masks. The two independent components
of $V_{0,t}$ have common face variances $1-t$ and $v/2$; at least one
is at least $1/3$. The retained summands have face norm at most
$2\delta$, and the expected face norms of the two components are
bounded by $2\E\|W\|_\infty$ and $g_\Sigma$, respectively.
Applying Lemma~\ref{lem:concentration} to the component with variance
at least $1/3$, with the other as an independent shift, gives
\eqref{hgap:eq:ac} with coefficient $\kappa$, uniformly in $t$.
This argument requires no lower bound on the pairwise gaps.
Also, $\delta\sqrt L\le\sigma_t<1$ once $\mathsf A\ge\sqrt2$.
The normalization $Y_t=\sqrt{1+\varepsilon}\,U_t$ changes the
smoothing width to $\widetilde h=\sqrt{1+\varepsilon}\,h$;
since $\widetilde h\asymp h$, we express the estimates below in
terms of $h$.

For signed face indices $j,k$, define
$\rho_{jk}:=u_j\cdot u_k$ and
$\nu_{jk}:=\sqrt{1-\rho_{jk}^2}$.
The Hessian decomposition~\eqref{cap:hessian} separates the
fourth-order remainder into single-face and unordered-pair
contributions. Group them as
\[
 R_{4,a,h}(t)=R_{\mathrm{sep},a,h}(t)+R_{\mathrm{exc},a,h}(t),
\]
where $R_{\mathrm{exc},a,h}$ contains the pair terms with
$0<1-|\rho_{jk}|<c$, and $R_{\mathrm{sep},a,h}$ contains all
single-face terms and the pair terms with $1-|\rho_{jk}|\ge c$.
There is no missing contribution at zero gap: if $u_k=\pm u_j$,
then the directional coordinates satisfy $z_k=\pm z_j$ and
$w_k=\pm w_j$, so the coefficient
$B_{jk}(z,w)$ in~\eqref{cap:pair-coefficient} is zero.

We first bound the exceptional contribution. For each such face
pair, the direct estimate~\eqref{cap:pair} gives
\[
 |R_{4,jk,a,h}(t)|\le\frac{C\delta^2\nu_{jk}}{h^2}.
\]
For two distinct original coordinates, all four choices of signs
give the same absolute correlation. If their coordinate gap is
$\eta$, then
$\nu_{jk}=\sqrt{\eta(2-\eta)}\le\sqrt{2\eta}$.
The two faces of a single coordinate have zero gap and contribute
zero. Thus each term of $B_c$ accounts for at most four unordered
signed-face pairs, and summation gives
\begin{equation}
 |R_{\mathrm{exc},a,h}(t)|
 \le\frac{C\delta^2B_c}{h^2},
 \label{budget:eq:exceptional}
\end{equation}
uniformly in $a$ and $t$. In particular, this bound stays uniform
as any positive exceptional gap tends to zero.

We next bound the separated contribution without imposing a gap
condition on the exceptional pairs. Set $d_c:=c(2-c)$, so that
$c\le d_c\le1$. Proposition~\ref{hgap:prop:cubic} has no gap
assumption. The angle-weighted bound in
Proposition~\ref{hgap:prop:pair} also holds before a minimum gap
is imposed, and its left side is a sum of nonnegative nuclear
norms. We may therefore discard the exceptional terms from that
sum. On each remaining pair,
\[
 \lambda_{jk}=\vartheta_{jk}\sin\vartheta_{jk}
 \ge\sin^2\vartheta_{jk}
 =1-\rho_{jk}^2\ge d_c,
 \qquad \vartheta_{jk}:=\arccos\rho_{jk}.
\]
Removing the angle weights from the restricted sum costs at most
$d_c^{-1}$. Only the indices in the pair sum have been restricted:
the smoothing function and its logistic weights still involve
all $m$ faces. Consequently the matrix estimates apply to the
same concentrated base verified above. Combining the two orders
of each pair in~\eqref{hgap:eq:hessian} gives the unordered
coefficient in~\eqref{cap:pair-coefficient}, so this restriction
selects exactly the pair terms in $R_{\mathrm{sep},a,h}$.

To pass from these matrices to the remainder, we use the
translated-weight expansion in the proof of
Lemma~\ref{hgap:thm:fourth}. Here are the hypotheses that make this
step valid. Let $\mathcal H$ be generated by the original summands,
masks, and $G_0$. In each term of
\eqref{hgap:eq:explicit-remainder}, the repeated direction $Z_i$
is a copy of $\zeta_i^0$. The array of these copies is independent
of $\mathcal H$ and satisfies
\[
 \sum_i\E[Z_iZ_i^\top\mid\mathcal H]=I_r.
\]
The other two directions have face norm at most $\delta$, and
the shift has face norm at most $4\delta$. The copies and shifts
are jointly independent of $G_1$ conditional on $\mathcal H$,
although a shift may depend on the copies. Expanding the translated
weights before taking absolute values separates the unshifted
Gaussian matrices from these copy-dependent coefficients. The
displayed covariance identity bounds the coefficient matrices,
and $h\ge\mathsf A\delta$ makes the translation series summable,
exactly as in Lemma~\ref{hgap:thm:fourth}. The concentration estimate
is used for the unconditional law of $V_{0,t}$; no concentration
conditional on $\mathcal H$ is asserted or needed.

The single-face matrix sum contributes $h\kappa L^{3/2}$, and
the restricted pair matrix sum contributes
$h^2\kappa L^{3/2}/(d_c\sigma_t)$. Their derivative factors are
$h^{-1}\sigma_t^{-3}$ and $h^{-2}\sigma_t^{-2}$, respectively.
The bounded integrals of the scalar kernels in
\eqref{hgap:eq:explicit-remainder} change only the absolute
constant. Hence
\begin{align}
 |R_{\mathrm{sep},a,h}(t)|
 &\le C\delta^2\left\{
    \frac{h\kappa L^{3/2}}{h\sigma_t^3}
   +\frac{h^2\kappa L^{3/2}}{h^2\sigma_t^2d_c\sigma_t}
                    \right\}\notag\\
 &\le\frac{C\delta^2\kappa L^{3/2}}{d_cv^{3/2}},
 \label{budget:eq:separated}
\end{align}
where the last inequality uses $d_c\le1$ and $\sigma_t^2=v/2$.

Near the start of the interpolation, a different estimate is
more useful. Proposition~\ref{an:prop:derivative} gives
\[
 |R_{4,a,h}(t)|\le\frac{C\delta^2\kappa L}{hv}
\]
for the full remainder. We do not need this estimate separately
for either part of the decomposition. Indeed, combining it with
\eqref{budget:eq:exceptional}--\eqref{budget:eq:separated} and using
$\min\{x,y+z\}\le\min\{x,y\}+z$ for nonnegative $x,y,z$ yields
\[
 |R_{4,a,h}(t)|\le
 C\min\left\{\frac{\delta^2\kappa L}{hv},
              \frac{\delta^2\kappa L^{3/2}}{d_cv^{3/2}}\right\}
 +\frac{C\delta^2B_c}{h^2}.
\]
Set $y:=v/\varepsilon$ and use $\sqrt\varepsilon=h\sqrt L$.
Extending the upper integration limit for the nonnegative minimum
to infinity gives
\begin{align}
 \sup_a\int_0^1|R_{4,a,h}(t)|\,dt
 &\le\frac{C\delta^2\kappa L}{h}
   \int_1^\infty\min\{y^{-1},d_c^{-1}y^{-3/2}\}\,dy
   +\frac{C\delta^2B_c}{h^2}\notag\\
 &\le C\left\{\frac{\delta^2\kappa LJ_c}{h}
                  +\frac{\delta^2B_c}{h^2}\right\}.
 \label{budget:eq:integrated}
\end{align}
For the last step, the two integrands agree at $y=d_c^{-2}$.
The integral before that point is $2\log(1/d_c)$, and the integral
after it is $2$. Their sum is at most $2J_c$ because $d_c\ge c$.
Thus the pointwise factor $d_c^{-1}$ costs only a logarithm after
time integration, while the exceptional-pair cost remains additive.

The smoothing, initialization, and cubic estimates used in
Appendix~\ref{app:correlation-gap-proof} contribute at most
$C\kappa L(h+\delta+\delta^2/h)$. Since $h\ge\mathsf A\delta$
and $J_c\ge1$, adding them to~\eqref{budget:eq:integrated} gives
\begin{equation}
 \rh(W,Z)\le C\left\{
       \kappa Lh+\frac{\delta^2\kappa LJ_c}{h}
                    +\frac{\delta^2B_c}{h^2}\right\}.
 \label{budget:eq:balance}
\end{equation}
There are two balances with the smoothing cost $\kappa Lh$.
The separated-pair term suggests width $\delta\sqrt{J_c}$;
the exceptional-pair term suggests width
$(\delta^2B_c/(\kappa L))^{1/3}$. Choose the larger width:
\[
 h:=\mathsf A\max\left\{\delta\sqrt{J_c},
                    \left(\frac{\delta^2B_c}{\kappa L}\right)^{1/3}
              \right\}.
\]
This gives $h\ge\mathsf A\delta$ and
\[
 \frac{\delta^2\kappa LJ_c}{h}
       \le\frac{\kappa Lh}{\mathsf A^2},\qquad
 \frac{\delta^2B_c}{h^2}
       \le\frac{\kappa Lh}{\mathsf A^3}.
\]
If $Lh^2<1/4$, substitution in~\eqref{budget:eq:balance}, followed by
\[
 \kappa Lh\le\mathsf A\left\{
       \delta\kappa L\sqrt{J_c}
       +\delta^{2/3}(\kappa L)^{2/3}B_c^{1/3}\right\},
\]
proves the desired bound with $\kappa$ in place of $g_\Sigma$.
If instead $Lh^2\ge1/4$, the expression in braces is at least
\[
 \frac{\kappa Lh}{\mathsf A}
 \ge\frac{\kappa\sqrt L}{2\mathsf A}
 \ge\frac1{2\mathsf A}.
\]
The trivial bound $\rh(W,Z)\le1$ therefore gives the same conclusion,
with the cap at $1$, even when this choice of width is outside the
range used to derive~\eqref{budget:eq:balance}.

It remains to replace $\kappa$ by $g_\Sigma$. Lemma~\ref{intrinsic:lem:maximum}
gives $\E\|W\|_\infty\le C(g_\Sigma+\delta L)$ independently of
the rectangle comparison being proved. If $\delta L\le g_\Sigma$,
this estimate and $g_\Sigma\ge\sqrt{2/\pi}$ imply
$\kappa\le Cg_\Sigma$. Substitution in both terms proves
\eqref{budget:eq:result}. If $\delta L>g_\Sigma$, the first term
in that bound already satisfies
\[
 \delta g_\Sigma L\sqrt{J_c}\ge g_\Sigma^2\ge2/\pi,
\]
so the trivial probability bound suffices. The argument includes
$B_c=0$ and singular covariance matrices, since zero-gap pairs
contribute zero and all derivative estimates were taken on the
Gaussian support.
\end{proof}
\clearpage
\begin{thebibliography}{9}
\setlength{\itemsep}{1pt}
\raggedright
\phantomsection
\addcontentsline{toc}{section}{References}
\bibitem{MainPaper2026}
V. Chernozhukov, D. Chetverikov, K. Kato, and Y. Koike (2026).
Sharp High-Dimensional Gaussian Approximations over Rectangles
without Covariance Restrictions.
16 November 2026.
\url{https://github.com/VC2015/HDCLTs/blob/main/research-rectangle-clt-singular/Rectangle_CLT_arXiv_2026-11-16.pdf}.

\bibitem{Prekopa1973}
A.~Pr\'ekopa (1973).
On logarithmic concave measures and functions.
\emph{Acta Scientiarum Mathematicarum (Szeged)} \textbf{34}, 335--343.
\url{https://rutcor.rutgers.edu/Prekopa/pdf/SCIENT2.pdf}.

\end{thebibliography}
\end{document}
